% concatenated from manuscript.tex
\documentclass[11pt,reqno]{amsart}

\usepackage[T1]{fontenc}
\usepackage[utf8]{inputenc}
\usepackage[expansion=false,protrusion=true]{microtype}
\usepackage{amsmath,amssymb,mathtools}
\usepackage{booktabs}
\usepackage{enumitem}
\usepackage{graphicx}
\usepackage[margin=2.6cm]{geometry}
\usepackage[colorlinks=true,linkcolor=black,citecolor=black,urlcolor=black]{hyperref}

\numberwithin{equation}{section}

\theoremstyle{plain}
\newtheorem{theorem}{Theorem}[section]
\newtheorem{proposition}[theorem]{Proposition}
\newtheorem{lemma}[theorem]{Lemma}
\newtheorem{corollary}[theorem]{Corollary}
\theoremstyle{definition}
\newtheorem{assumption}[theorem]{Assumption}
\theoremstyle{remark}
\newtheorem{remark}[theorem]{Remark}

\newcommand{\bT}{\mathbf{T}}
\newcommand{\bZ}{\mathbf{Z}}
\newcommand{\bU}{\mathbf{U}}
\newcommand{\bV}{\mathbf{V}}
\newcommand{\bW}{\mathbf{W}}
\newcommand{\bB}{\mathbf{B}}
\newcommand{\bDelta}{\boldsymbol{\Delta}}
\newcommand{\bone}{\mathbf{1}}
\newcommand{\ba}{\mathbf{a}}
\newcommand{\bze}{\mathbf{0}}
\newcommand{\GH}{\Gamma_H}
\newcommand{\Hmax}{H_{\max}}
\newcommand{\Hmin}{H_{\min}}
\newcommand{\R}{\mathbb{R}}
\newcommand{\N}{\mathbb{N}}
\newcommand{\Prob}{\mathbb{P}}
\newcommand{\E}{\mathbb{E}}
\newcommand{\Norm}{\mathcal{N}}
\newcommand{\Hi}{\mathfrak{H}}
\DeclareMathOperator{\Var}{Var}
\DeclareMathOperator{\Cov}{Cov}
\DeclareMathOperator{\Bias}{Bias}
\DeclareMathOperator{\tr}{tr}
\DeclareMathOperator{\rank}{rank}
\DeclareMathOperator{\MSE}{MSE}

\emergencystretch=2.5em

\begin{document}

\title[Exact inference for multi-mixed fractional Brownian motion]
      {Exact finite-sample inference for multi-mixed\\
       fractional Brownian motion with drift}

\author{Afrah Al-Harby}
\address{Department of Mathematics, College of Science, Imam Abdulrahman Bin Faisal
University, P.\,O.\ Box 1982, Dammam, Saudi Arabia}
\email{2250500238@iau.edu.sa}

\author{Ezzedine Mliki}
\address{Department of Mathematics, College of Science, Imam Abdulrahman Bin Faisal
University, P.\,O.\ Box 1982, Dammam, Saudi Arabia}
\address{Basic and Applied Scientific Research Center, Imam Abdulrahman Bin Faisal
University, P.\,O.\ Box 1982, 31441, Dammam, Saudi Arabia}
\email{ermliki@iau.edu.sa}

\begin{abstract}
In this paper we study a linear drift perturbed by a superposition of $m$ independent
fractional Brownian motions with known Hurst parameters and a common scale, observed at
$N$ equidistant times. Inference for such models is usually asymptotic; we show that here
it is exact. We derive the maximum likelihood estimators of the drift $\theta$ and of the
scale $\alpha^{2}$ in closed form and obtain their exact finite-sample joint law:
$\widehat\theta$ is Gaussian, $N\widehat\alpha^{\,2}/\alpha^{2}$ is chi-square with $N-1$
degrees of freedom, and the two are independent. As this law is free of every model
parameter, we deduce Student and chi-square confidence intervals and tests of exact level
for every $N\ge2$, whatever the Hurst vector. We also prove that the estimators are
uniformly minimum variance unbiased with $\widehat\theta$ attaining the Cram\'er--Rao
bound at every $N$, that both are strongly consistent and asymptotically normal, and that
the drift estimators form, in law, a Brownian motion run along their own variance scale.
A sharp non-asymptotic bound shows that the accuracy of the drift is governed by the
length of the observation window and not by the mesh, and a Monte Carlo study confirms
exact coverage, even at small sample sizes, and quantifies what is lost when the Hurst
vector is misspecified.
\end{abstract}

\keywords{Multi-mixed fractional Brownian motion; maximum likelihood estimation;
discrete observations; exact finite-sample distribution; long-range dependence; strong
consistency; asymptotic normality}

\subjclass[2020]{Primary 62M09; Secondary 60G22, 60G15, 62F12, 62F25, 60H07}

\maketitle

\section{Introduction}

Fractional Brownian motion (fBm) is the centred continuous-time Gaussian process
$B^{H}=\{B^{H}(t),\,t\ge0\}$ introduced by Mandelbrot and Van Ness \cite{MVN68}, indexed
by a Hurst parameter $H\in(0,1)$ and characterised by $B^{H}(0)=0$ and the covariance
function
\begin{equation}\label{eq:fbmcov}
  \E\bigl[B^{H}(t)B^{H}(s)\bigr]=\tfrac12\bigl(t^{2H}+s^{2H}-|t-s|^{2H}\bigr),
  \qquad s,t\ge0 .
\end{equation}
The Hurst parameter governs both the dependence structure and the regularity of the
sample paths. The value $H=\frac12$ corresponds to standard Brownian motion,
$H>\frac12$ yields positively correlated increments and long-range dependence, while
$H<\frac12$ produces negatively correlated (antipersistent) increments. A key feature of
fBm is the stationarity of its increments,
\begin{equation}\label{eq:statincr}
  B^{H}(t+s)-B^{H}(t)\overset{d}{=}B^{H}(s),\qquad s,t\ge0,
\end{equation}
which plays a central role in the statistical analysis of discretely observed processes,
where inference is naturally formulated in terms of the increment process. Owing to its
ability to capture simultaneously self-similarity and long-range dependence, fBm has
become a standard model for complex temporal phenomena in hydrology, geosciences,
telecommunication networks, economics and finance. In mathematical finance, fractional
and mixed fractional models have been used in option pricing to incorporate long memory
and departures from the behaviour predicted by classical Brownian models; see, for
instance, Haddadi and Nasrollahi \cite{HN25}, who combine fractional and mixed fractional
models with Gram--Charlier expansions for European options under non-normal returns.

Despite its versatility, a single fractional component may be insufficient for phenomena
exhibiting several scaling regimes or several distinct sources of temporal dependence.
This limitation motivates multi-mixed fractional models obtained by superimposing
finitely many independent fractional Brownian motions with distinct Hurst parameters.
Such models offer greater flexibility and a richer covariance structure while remaining
Gaussian, hence amenable to likelihood-based inference. The terminology follows Maleki
Almani and Sottinen \cite{MS23}, who study the path properties of such superpositions.
The structural theory of mixtures of this kind is by now well developed. Alajmi and Mliki
\cite{AM21} introduce the mixed generalized fractional Brownian motion, which unifies the
mixed fractional Brownian motion of Cheridito and the generalized fractional Brownian
motion of Zili, and establish its self-similarity and its long-range dependence together
with conditions under which it fails to be a semimartingale. What such results do not
provide is inference: the sampling behaviour of the parameters of these processes, when
the process is observed at finitely many times, is a separate question and the one
addressed here.

Parameter estimation for fractional Gaussian models has received considerable attention.
For the drifted fractional Brownian motion $Y(t)=\mu t+\sigma B^{H}(t)$ observed at
discrete times, Hu, Nualart, Xiao and Zhang \cite{HNXZ11} developed an exact
likelihood-based framework, deriving closed-form maximum likelihood estimators for the
drift and variance parameters together with their asymptotic properties. This framework
has since been extended in several directions. Mishura, Ralchenko and Shklyar
\cite{MRS17} generalised the likelihood approach to Gaussian processes with stationary
increments, which contains fBm as a particular case. Xiao, Zhang and Zhang \cite{XZZ11}
studied maximum likelihood estimation for models driven by mixed fractional Brownian
motion, that is, by the sum of an independent standard Brownian motion and a fractional
Brownian motion. Mishura and Voronov \cite{MV15} considered maximum likelihood estimation
of the drift for a model driven by the sum of two independent fractional Brownian motions
with distinct Hurst parameters, reducing the problem to a Fredholm integral equation of
the second kind.

The closest current work allows the components to carry \emph{separate} scales.
Ralchenko and Yakovliev study $X_t=\theta t+\sigma W_t+\kappa B^{H}_t$ in \cite{RY23} and
$X_t=\kappa B^{H_1}_t+\sigma B^{H_2}_t$ in \cite{RY24}, constructing strongly consistent
estimators of all the parameters and proving joint asymptotic normality; the price of the
second free scale is that these limit theorems require $H<\tfrac12$ in the first model
and $H_1<H_2<\tfrac34$ in the second. For the sub-fractional model with drift, Kuang and
Liu \cite{KL15} obtain the central limit theorem together with Berry--Ess\'een bounds,
which are needed precisely because the finite-sample law is not available. Every result
in this line is asymptotic. The present paper takes the opposite trade: a single scale
shared by all components, and in exchange an exact finite-sample theory that holds at
every $N\ge2$ and imposes no restriction on the Hurst vector.

\subsection*{The model and our contribution}
Let $t_k=kh$, $k=1,\dots,N$, with mesh $h>0$, and consider the observation vector
\begin{equation}\label{eq:obsvec}
  \bZ=\bigl(Z(t_1),\dots,Z(t_N)\bigr)^{\top}=\theta\bT+\alpha\sum_{r=1}^{m}\bB^{H_r},
\end{equation}
generated by the continuous-time model
\begin{equation}\label{eq:model}
  Z(t)=\theta t+\alpha\sum_{r=1}^{m}B^{H_r}(t),\qquad t\ge0,
\end{equation}
where $B^{H_1},\dots,B^{H_m}$ are independent fractional Brownian motions with known
Hurst parameters $H_1,\dots,H_m\in(1/2,1)$, and
\[
  \bT=(t_1,\dots,t_N)^{\top},\qquad
  \bB^{H_r}=\bigl(B^{H_r}(t_1),\dots,B^{H_r}(t_N)\bigr)^{\top}.
\]
Under the independence assumption, $\bZ$ is a Gaussian vector with mean $\theta\bT$ and
covariance matrix $\alpha^{2}\GH$, where $\GH=\sum_{r=1}^{m}\Gamma_{H_r}$ and
\begin{equation}\label{eq:GammaHr}
  (\Gamma_{H_r})_{ij}=\tfrac12\bigl(t_i^{2H_r}+t_j^{2H_r}-|t_i-t_j|^{2H_r}\bigr),
  \qquad i,j=1,\dots,N .
\end{equation}
Model \eqref{eq:model} is the common-scale specialisation of the general multi-mixed
model $Z(t)=\theta t+\sum_{r=1}^{m}\alpha_rB^{H_r}(t)$. The restriction is natural
whenever the several sources of dependence are believed to act at a common intensity and
to differ only in memory, as when one mechanism is observed through channels of differing
persistence; it is also what the classical mixed models reduce to after
reparameterisation, since $\sigma W+\kappa B^{H}$ may be written as
$\sigma(W+\rho B^{H})$ with $\rho=\kappa/\sigma$ fixed. Analytically the common scale is
what makes the likelihood explicitly solvable, because the covariance matrix is then a
\emph{known} matrix multiplied by a single unknown constant; and the resulting inference
problem already exhibits the multi-scale covariance structure that distinguishes the
multi-mixed setting from the single-component one.

Our contribution is the following.
\begin{enumerate}[label=\textup{(\arabic*)}]
\item We derive closed-form maximum likelihood estimators $\widehat\theta$ and
      $\widehat\alpha^{\,2}$ (Theorem \ref{th:mle}), and we show that the whole
      multi-mixed structure enters the inference problem through a single scalar, the
      information quantity $d_N=\bT^{\top}\GH^{-1}\bT$, which admits an exact
      representation in terms of the increment process (Proposition \ref{pr:increment}).
\item We obtain the exact finite-sample joint law of
      $(\widehat\theta,\widehat\alpha^{\,2})$: the two estimators are independent,
      $\widehat\theta$ is Gaussian, and $N\widehat\alpha^{\,2}/\alpha^{2}$ is chi-square
      with $N-1$ degrees of freedom (Theorem \ref{th:exactlaw}). All moment formulas
      follow at once (Corollary \ref{co:moments}), and exact confidence intervals and
      exact tests for $\theta$ and $\alpha^{2}$ become available
      (Corollaries \ref{co:ci} and \ref{co:tests}).
\item We establish the optimality of the estimators in the finite-sample sense: the pair
      $(\widehat\theta,\widehat\alpha^{\,2})$ is complete and sufficient,
      $\widehat\theta$ and the bias-corrected estimator
      $\widetilde\alpha^{\,2}=N\widehat\alpha^{\,2}/(N-1)$ are uniformly minimum variance
      unbiased, and $\widehat\theta$ attains the Cram\'er--Rao bound at every sample size
      (Theorem \ref{th:umvu} and Proposition \ref{pr:fisher}).
\item We establish an explicit non-asymptotic variance bound for $\widehat\theta$
      (Proposition \ref{pr:varbound}) showing that the accuracy of drift estimation is
      driven by the length $t_N=Nh$ of the observation window and by
      $\Hmax=\max_r H_r$, and we show that the order of this bound is sharp
      (Remark \ref{re:sharp}).
\item We prove strong consistency of both estimators (Theorems \ref{th:sctheta} and
      \ref{th:scalpha}) and asymptotic normality (Theorem \ref{th:an}), the latter both
      by an elementary argument and by the Malliavin calculus criterion of Nualart and
      Ortiz-Latorre \cite{NOL08}.
\item We show that the sequence $\{\widehat\theta^{(N)}\}_{N\ge1}$ has independent
      increments and is, in law, a Brownian motion evaluated along the decreasing
      sequence of its own variances (Theorem \ref{th:indincr} and
      Corollary \ref{co:tcbm}); this yields a second, structural proof of strong
      consistency.
\item A Monte Carlo study quantifies the finite-sample behaviour of the estimators and
      confirms the exact distributional results, including the coverage of the exact
      confidence intervals (Section \ref{se:sim}). A final experiment measures the cost
      of a misspecified Hurst vector and shows that the interval for $\theta$ degrades
      slowly, and in a direction one can predict, whereas the interval for
      $\alpha^{2}$ does not (Section \ref{sse:misspec}).
\end{enumerate}

Throughout, the asymptotic regime is the increasing domain one: the mesh $h>0$ is held
fixed and $N\to\infty$, so that the observation window $t_N=Nh$ grows without bound.

\section{Model specification and likelihood framework}

\subsection{Notation and standing assumptions}
Let $(\Omega,\mathcal F,\Prob)$ be a complete probability space carrying independent
fractional Brownian motions $\{B^{H_r}(t),\,t\ge0\}$, $r=1,\dots,m$, and write
$H=(H_1,\dots,H_m)^{\top}$ for the vector of Hurst parameters.

\begin{assumption}\label{as:main}
The mesh $h>0$ and the number of components $m\in\N$ are fixed and known; the Hurst
vector $H$ is known with $H_r\in(1/2,1)$ for every $r$; the drift $\theta\in\R$ and the
scale $\alpha>0$ are unknown. The observation times are $t_k=kh$, $k=1,\dots,N$, and we
set $t_0=0$.
\end{assumption}

\begin{remark}\label{re:half}
The restriction $H_r>1/2$ reflects the long-memory motivation of the model and is not
needed for most of the analysis: Theorems \ref{th:mle}, \ref{th:exactlaw},
\ref{th:umvu} and \ref{th:indincr}, Propositions \ref{pr:increment} and \ref{pr:fisher}
and Corollaries \ref{co:ci} and \ref{co:tests} hold verbatim for arbitrary
$H_r\in(0,1)$, and the asymptotic statements of Section \ref{se:asymp} hold for arbitrary
$H_r\in(0,1)$ with $\Hmax=\max_rH_r$ throughout. This is what allows the degenerate cases
listed in Remark \ref{re:special} to be covered. Only Remark \ref{re:sharp}, which
identifies the exact order of the variance of $\widehat\theta$, uses $\Hmax>1/2$.
\end{remark}

Define the superposition
\begin{equation}\label{eq:U}
  U(t)=\sum_{r=1}^{m}B^{H_r}(t),\qquad t\ge0,
\end{equation}
which we call a multi-mixed fractional Brownian motion, and the associated observation
vectors
\[
  \bU=\bigl(U(t_1),\dots,U(t_N)\bigr)^{\top}=\sum_{r=1}^{m}\bB^{H_r},
  \qquad \bT=(t_1,\dots,t_N)^{\top}.
\]
The observation vector is $\bZ=\theta\bT+\alpha\bU$. By independence of the components,
\begin{equation}\label{eq:Zlaw}
  \bZ\sim\Norm\bigl(\theta\bT,\ \alpha^{2}\GH\bigr),\qquad
  \GH=\sum_{r=1}^{m}\Gamma_{H_r},
\end{equation}
with $\Gamma_{H_r}$ as in \eqref{eq:GammaHr}. We write $\GH$ for $\GH^{(N)}$ whenever the
dependence on $N$ is clear, and $\Gamma_{H,N}$ when it must be displayed.

\begin{lemma}\label{le:nondeg}
Under Assumption \ref{as:main}, $\GH$ is symmetric positive definite; in particular
$\GH^{-1}$ exists and $\bT^{\top}\GH^{-1}\bT>0$.
\end{lemma}

\begin{proof}
Fix $r$ and set $H=H_r$. Let $\Delta^{B}_k=B^{H}(t_k)-B^{H}(t_{k-1})$, $k=1,\dots,N$, so
that $\bB^{H}=A\Delta^{B}$ with $\Delta^{B}=(\Delta^{B}_1,\dots,\Delta^{B}_N)^{\top}$ and
$A=(\mathbf 1_{\{j\le i\}})_{i,j=1}^{N}$ lower triangular with unit diagonal, hence
invertible. Because the grid is equidistant and $B^{H}$ has stationary increments
(see \eqref{eq:statincr}), $\Delta^{B}$ is a fractional Gaussian noise of step $h$, that
is, a stationary Gaussian sequence whose spectral density is strictly positive almost
everywhere on $[-\pi,\pi]$ \cite[Ch.~1]{Mis08}. A stationary sequence with an almost
everywhere positive spectral density has nonsingular covariance matrices at every order,
so $C^{B}_N:=\Cov(\Delta^{B})$ is positive definite, and therefore
$\Gamma_{H}=AC^{B}_NA^{\top}$ is positive definite as well. Summing over $r$,
$\GH=\sum_r\Gamma_{H_r}$ is positive definite. Finally $\bT\ne\bze$ because
$t_N=Nh>0$, whence $\bT^{\top}\GH^{-1}\bT>0$.
\end{proof}

\subsection{The increment process}

\begin{lemma}\label{le:statincr}
The process $U$ defined in \eqref{eq:U} is a centred Gaussian process with $U(0)=0$ and
stationary increments: for every $s\ge0$ the shifted increment process
$\{U(t+s)-U(s),\,t\ge0\}$ has the same law as $\{U(t),\,t\ge0\}$. Moreover, for all
$s,t\ge0$,
\begin{equation}\label{eq:Ucov}
  \E\bigl[U(t)U(s)\bigr]
  =\tfrac12\sum_{r=1}^{m}\bigl(t^{2H_r}+s^{2H_r}-|t-s|^{2H_r}\bigr),
\end{equation}
and the increment vector $\bDelta=(\Delta_1,\dots,\Delta_N)^{\top}$,
$\Delta_k=U(t_k)-U(t_{k-1})$, has the symmetric Toeplitz covariance matrix
$C_N=\bigl(c(|i-j|)\bigr)_{i,j=1}^{N}$ with
\begin{equation}\label{eq:CN}
  c(k)=\tfrac12\sum_{r=1}^{m}h^{2H_r}
       \bigl(|k+1|^{2H_r}+|k-1|^{2H_r}-2|k|^{2H_r}\bigr),\qquad k\ge0 .
\end{equation}
\end{lemma}

\begin{proof}
$U$ is a finite sum of independent centred Gaussian processes vanishing at the origin,
hence centred Gaussian with $U(0)=0$; formula \eqref{eq:Ucov} follows from
\eqref{eq:fbmcov} and independence.

Both $\{U(t+s)-U(s)\}_{t\ge0}$ and $\{U(t)\}_{t\ge0}$ are centred Gaussian, so it
suffices to compare their covariance functions. Fix $s\ge0$ and $0\le u\le v$. Using
\eqref{eq:Ucov},
\begin{align*}
  \E\bigl[(U(u+s)-U(s))(U(v+s)-U(s))\bigr]
  &=\E[U(u+s)U(v+s)]-\E[U(u+s)U(s)]\\
  &\qquad-\E[U(s)U(v+s)]+\E[U(s)^{2}]\\
  &=\tfrac12\sum_{r=1}^{m}\Bigl[(u+s)^{2H_r}+(v+s)^{2H_r}-|u-v|^{2H_r}\Bigr]\\
  &\quad-\tfrac12\sum_{r=1}^{m}\Bigl[(u+s)^{2H_r}+s^{2H_r}-u^{2H_r}\Bigr]\\
  &\quad-\tfrac12\sum_{r=1}^{m}\Bigl[s^{2H_r}+(v+s)^{2H_r}-v^{2H_r}\Bigr]
        +\sum_{r=1}^{m}s^{2H_r}\\
  &=\tfrac12\sum_{r=1}^{m}\bigl(u^{2H_r}+v^{2H_r}-|u-v|^{2H_r}\bigr)
   =\E\bigl[U(u)U(v)\bigr],
\end{align*}
all terms in $(u+s)^{2H_r}$, $(v+s)^{2H_r}$ and $s^{2H_r}$ cancelling pairwise. Since a
centred Gaussian process is determined by its covariance function, the two processes have
the same law.

Stationarity of the increments and the equidistance of the grid imply that
$\E[\Delta_i\Delta_j]$ depends on $i,j$ only through $|i-j|$, and taking $t=|i-j|h$,
$s=h$ in \eqref{eq:Ucov} gives
$\E[\Delta_i\Delta_j]=\tfrac12\sum_r\bigl((|i-j|+1)^{2H_r}h^{2H_r}
+(|i-j|-1)^{2H_r}h^{2H_r}-2|i-j|^{2H_r}h^{2H_r}\bigr)$, which is \eqref{eq:CN}.
\end{proof}

\begin{remark}\label{re:notss}
Unlike each of its components, $U$ is not self-similar when the $H_r$ are not all equal:
by \eqref{eq:Ucov}, $\Var U(t)=\sum_rt^{2H_r}$ is not a single power of $t$. It is
exactly this superposition of scaling exponents that gives the model its multi-scale
character, and it is why the aggregated matrix $\GH=\sum_r\Gamma_{H_r}$, rather than any
single $\Gamma_{H_r}$, must be inverted in the likelihood.
\end{remark}

\subsection{The likelihood function}

\begin{theorem}\label{th:loglik}
Under Assumption \ref{as:main}, the log-likelihood of $(\theta,\alpha^{2})$ based on
$\bZ$ is
\begin{equation}\label{eq:loglik}
  \ell(\theta,\alpha^{2};\bZ)=-\frac N2\log(2\pi)-\frac N2\log\alpha^{2}
  -\frac12\log|\GH|-\frac{1}{2\alpha^{2}}(\bZ-\theta\bT)^{\top}\GH^{-1}(\bZ-\theta\bT).
\end{equation}
\end{theorem}

\begin{proof}
By \eqref{eq:Zlaw} and Lemma \ref{le:nondeg}, $\bZ$ admits the nondegenerate Gaussian
density
\[
  f(\bZ;\theta,\alpha^{2})=(2\pi\alpha^{2})^{-N/2}|\GH|^{-1/2}
  \exp\Bigl\{-\tfrac{1}{2\alpha^{2}}(\bZ-\theta\bT)^{\top}\GH^{-1}(\bZ-\theta\bT)\Bigr\},
\]
where we used $|\alpha^{2}\GH|=(\alpha^{2})^{N}|\GH|$. Taking logarithms gives
\eqref{eq:loglik}.
\end{proof}

\section{Maximum likelihood estimators and their exact distribution}\label{se:mle}

\subsection{The estimators}

\begin{theorem}\label{th:mle}
The log-likelihood \eqref{eq:loglik} admits the unique global maximiser
\begin{equation}\label{eq:mle}
  \widehat\theta=\frac{\bT^{\top}\GH^{-1}\bZ}{\bT^{\top}\GH^{-1}\bT},
  \qquad
  \widehat\alpha^{\,2}=\frac1N\left(\bZ^{\top}\GH^{-1}\bZ
    -\frac{(\bT^{\top}\GH^{-1}\bZ)^{2}}{\bT^{\top}\GH^{-1}\bT}\right),
\end{equation}
on $\R\times(0,\infty)$, almost surely.
\end{theorem}

\begin{proof}
Write $s=\alpha^{2}>0$ and $Q(\theta)=(\bZ-\theta\bT)^{\top}\GH^{-1}(\bZ-\theta\bT)$.
For fixed $s$, $\ell$ is, up to an additive constant, equal to $-Q(\theta)/(2s)$. Since
$\GH^{-1}$ is positive definite, $Q$ is a strictly convex quadratic in $\theta$ with
\[
  Q'(\theta)=-2\,\bT^{\top}\GH^{-1}(\bZ-\theta\bT),\qquad
  Q''(\theta)=2\,\bT^{\top}\GH^{-1}\bT>0,
\]
so $Q$ has the unique minimiser $\widehat\theta$ given in \eqref{eq:mle}, for every
$s>0$. Substituting yields the profile log-likelihood
\[
  \ell_p(s)=-\frac N2\log(2\pi)-\frac N2\log s-\frac12\log|\GH|
            -\frac{Q(\widehat\theta)}{2s},
\]
whose derivative $\ell_p'(s)=-N/(2s)+Q(\widehat\theta)/(2s^{2})$ vanishes only at
$s=Q(\widehat\theta)/N$ and changes sign there from $+$ to $-$; hence this point is the
unique global maximiser. Note $Q(\widehat\theta)>0$ almost surely, because
$Q(\widehat\theta)=0$ would force $\bZ$ to be collinear with $\bT$, an event of
probability zero by Lemma \ref{le:nondeg}. Finally, expanding
\[
  Q(\widehat\theta)=\bZ^{\top}\GH^{-1}\bZ-2\widehat\theta\,\bT^{\top}\GH^{-1}\bZ
  +\widehat\theta^{\,2}\bT^{\top}\GH^{-1}\bT
  =\bZ^{\top}\GH^{-1}\bZ-\frac{(\bT^{\top}\GH^{-1}\bZ)^{2}}{\bT^{\top}\GH^{-1}\bT}
\]
gives the stated form of $\widehat\alpha^{\,2}=Q(\widehat\theta)/N$.
\end{proof}

\begin{remark}\label{re:gls}
$\widehat\theta$ is the generalised least squares estimator of the drift, which coincides
with the MLE under Gaussianity. Both estimators are computable from the data without
knowledge of $\alpha$, and $\widehat\theta$ does not depend on $\alpha$ at all.
\end{remark}

Throughout what follows we write
\begin{equation}\label{eq:dN}
  d_N:=\bT^{\top}\GH^{-1}\bT>0 ,
\end{equation}
the Fisher information about $\theta$ per unit of $\alpha^{2}$; see Proposition
\ref{pr:fisher}. The next result identifies $d_N$ in terms of the increment process and
shows that the multi-mixed structure enters the problem only through the Toeplitz matrix
$C_N$ of \eqref{eq:CN}. It will be the basis of the sharpness statement of
Remark \ref{re:sharp}, and it is also the form in which $\widehat\theta$ is best
computed.

\begin{proposition}\label{pr:increment}
Let $\bDelta^{Z}$ be the vector of observed increments, with $k$-th entry
$Z(t_k)-Z(t_{k-1})$, and let $C_N$ be the Toeplitz matrix \eqref{eq:CN}. Then
\begin{equation}\label{eq:dNincr}
  d_N=h^{2}\,\bone^{\top}C_N^{-1}\bone,
  \qquad
  \widehat\theta=\frac{\bone^{\top}C_N^{-1}\bDelta^{Z}}
                      {h\,\bone^{\top}C_N^{-1}\bone},
\end{equation}
where $\bone=(1,\dots,1)^{\top}\in\R^{N}$. Equivalently, $h\widehat\theta$ is the
generalised least squares estimator of the common mean $\E[\Delta^{Z}_k]=\theta h$ of the
stationary Gaussian sequence $\bDelta^{Z}$.
\end{proposition}

\begin{proof}
Let $A=(\mathbf 1_{\{j\le i\}})_{i,j=1}^{N}$, which is lower triangular with unit
diagonal and therefore invertible, with $A^{-1}$ the first-difference operator. Then
$\bZ=A\bDelta^{Z}$ and, because $t_k=kh$, $\bT=A(h\bone)$. Since
$\bDelta^{Z}-\theta h\bone=\alpha\bDelta$, Lemma \ref{le:statincr} gives
$\Cov(\bDelta^{Z})=\alpha^{2}C_N$, and therefore
$\alpha^{2}\GH=\Cov(\bZ)=A\Cov(\bDelta^{Z})A^{\top}=\alpha^{2}AC_NA^{\top}$, that is
$\GH=AC_NA^{\top}$ and $\GH^{-1}=A^{-\top}C_N^{-1}A^{-1}$. Substituting $A^{-1}\bT=h\bone$ and
$A^{-1}\bZ=\bDelta^{Z}$ in \eqref{eq:dN} and \eqref{eq:mle} gives
\[
  d_N=(A^{-1}\bT)^{\top}C_N^{-1}(A^{-1}\bT)=h^{2}\bone^{\top}C_N^{-1}\bone,
  \qquad
  \bT^{\top}\GH^{-1}\bZ=h\,\bone^{\top}C_N^{-1}\bDelta^{Z},
\]
which is \eqref{eq:dNincr}. The last assertion is the definition of the generalised least
squares estimator of the mean of a stationary sequence with covariance matrix
proportional to $C_N$.
\end{proof}

\subsection{The canonical reduction}
The whole distribution theory of this section rests on the following elementary change of
variables, which diagonalises the problem.

\begin{lemma}\label{le:canon}
Set
\begin{equation}\label{eq:canon}
  \bV=\GH^{-1/2}\bU,\qquad
  \ba=\frac{\GH^{-1/2}\bT}{\sqrt{d_N}},\qquad
  P=I_N-\ba\ba^{\top}.
\end{equation}
Then $\bV\sim\Norm(\bze,I_N)$, $\|\ba\|=1$, $P$ is the orthogonal projection onto
$\ba^{\perp}$ with $\rank P=N-1$, and
\begin{equation}\label{eq:canonid}
  \widehat\theta-\theta=\frac{\alpha\,\ba^{\top}\bV}{\sqrt{d_N}},
  \qquad
  \frac{N\widehat\alpha^{\,2}}{\alpha^{2}}=\bV^{\top}P\bV .
\end{equation}
\end{lemma}

\begin{proof}
By Lemma \ref{le:nondeg}, $\GH$ has a symmetric positive definite square root, and
$\bU\sim\Norm(\bze,\GH)$ gives $\bV=\GH^{-1/2}\bU\sim\Norm(\bze,I_N)$. Clearly
$\|\ba\|^{2}=\bT^{\top}\GH^{-1}\bT/d_N=1$, so $P=I_N-\ba\ba^{\top}$ satisfies
$P=P^{\top}=P^{2}$, $P\ba=\bze$ and $\tr P=N-1$.

Since $\bZ-\theta\bT=\alpha\bU$ and
$\bT^{\top}\GH^{-1}\bU=\bT^{\top}\GH^{-1/2}\bV=\sqrt{d_N}\,\ba^{\top}\bV$, we get from
\eqref{eq:mle}
\[
  \widehat\theta-\theta=\frac{\bT^{\top}\GH^{-1}(\bZ-\theta\bT)}{d_N}
  =\frac{\alpha\sqrt{d_N}\,\ba^{\top}\bV}{d_N}
  =\frac{\alpha\,\ba^{\top}\bV}{\sqrt{d_N}} .
\]
For the second identity, note that
$N\widehat\alpha^{\,2}=(\bZ-\widehat\theta\bT)^{\top}\GH^{-1}(\bZ-\widehat\theta\bT)
=\bigl\|\GH^{-1/2}(\bZ-\widehat\theta\bT)\bigr\|^{2}$ and
\[
  \GH^{-1/2}(\bZ-\widehat\theta\bT)=\alpha\bV-(\widehat\theta-\theta)\GH^{-1/2}\bT
  =\alpha\bV-\alpha(\ba^{\top}\bV)\ba=\alpha P\bV .
\]
Hence $N\widehat\alpha^{\,2}=\alpha^{2}\|P\bV\|^{2}=\alpha^{2}\bV^{\top}P\bV$.
\end{proof}

\begin{remark}\label{re:glm}
Lemma \ref{le:canon} says that the whitening map $z\mapsto\GH^{-1/2}z$ turns the
observation vector into
$\bV_{\!Z}:=\GH^{-1/2}\bZ\sim\Norm\bigl(\theta\sqrt{d_N}\,\ba,\ \alpha^{2}I_N\bigr)$
with $\ba$ a \emph{known} unit vector. After an orthogonal change of coordinates carrying
$\ba$ to the first axis, the model becomes the classical Gaussian linear model with a
single known regressor: one coordinate distributed as
$\Norm(\theta\sqrt{d_N},\alpha^{2})$ and $N-1$ coordinates i.i.d.\ $\Norm(0,\alpha^{2})$,
all independent. The exact theory of this section is therefore the exact small-sample
theory of that model transported through $\GH^{-1/2}$, and the entire multi-mixed
dependence structure (the number $m$ of components, the Hurst vector $H$ and the mesh
$h$) is compressed into the single scalar $d_N$. This is what makes exact inference
possible for arbitrary $m$ and arbitrary $H$, and it is also the reason why the results
of this section cannot be expected to survive the relaxation of the common-scale
assumption discussed in Section \ref{se:disc}: with component-specific scales
$\alpha_1,\dots,\alpha_m$ the covariance matrix
$\sum_r\alpha_r^{2}\Gamma_{H_r}$ is no longer a scalar multiple of a known matrix, and no
single whitening transformation removes the nuisance structure.
\end{remark}

\subsection{Exact finite-sample distribution}

\begin{theorem}\label{th:exactlaw}
Under Assumption \ref{as:main}, for every $N\ge2$:
\begin{enumerate}[label=\textup{(\roman*)}]
\item $\widehat\theta\sim\Norm\bigl(\theta,\ \alpha^{2}/d_N\bigr)$; in particular
      $\widehat\theta$ is unbiased;
\item $\dfrac{N\widehat\alpha^{\,2}}{\alpha^{2}}\sim\chi^{2}_{N-1}$;
\item $\widehat\theta$ and $\widehat\alpha^{\,2}$ are independent.
\end{enumerate}
\end{theorem}

\begin{proof}
We use Lemma \ref{le:canon}. Since $\bV\sim\Norm(\bze,I_N)$ and $\|\ba\|=1$, the scalar
$\ba^{\top}\bV$ is standard normal, and the first identity in \eqref{eq:canonid} gives
(i).

For (ii), $P$ is an orthogonal projection of rank $N-1$; writing
$P=O^{\top}\operatorname{diag}(I_{N-1},0)O$ for an orthogonal $O$ and using the
rotational invariance of $\Norm(\bze,I_N)$, we obtain
$\bV^{\top}P\bV=\sum_{i=1}^{N-1}\widetilde V_i^{2}\sim\chi^{2}_{N-1}$, where
$\widetilde{\bV}=O\bV\sim\Norm(\bze,I_N)$. The second identity in \eqref{eq:canonid}
gives (ii).

For (iii), the vector $(\ba^{\top}\bV,(P\bV)^{\top})^{\top}$ is Gaussian and
\[
  \Cov\bigl(P\bV,\ \ba^{\top}\bV\bigr)=P\,\E[\bV\bV^{\top}]\,\ba=P\ba=\bze,
\]
so $\ba^{\top}\bV$ and $P\bV$ are independent. As $\widehat\theta$ is a function of
$\ba^{\top}\bV$ and $\widehat\alpha^{\,2}=\alpha^{2}\|P\bV\|^{2}/N$ is a function of
$P\bV$, the two estimators are independent.
\end{proof}

\begin{remark}\label{re:pivot}
Statement (ii) is remarkable in that the law of $\widehat\alpha^{\,2}/\alpha^{2}$ depends
on the sample size only: it is free of the drift $\theta$, of the mesh $h$, of the number
of components $m$ and of the Hurst vector $H$. Thus the whole multi-scale dependence
structure of the model is absorbed by the whitening transformation $\bV=\GH^{-1/2}\bU$
and leaves no trace in the sampling distribution of the scale estimator. This is
confirmed numerically in Section \ref{se:sim}, where the standard deviation of
$\widehat\alpha$ is seen to be the same for $h=1/252$ and $h=1/12$. By contrast, the law
of $\widehat\theta$ does depend on $h$ and $H$, through the single scalar $d_N$.
\end{remark}

\begin{corollary}\label{co:moments}
For every $N\ge2$,
\begin{equation}\label{eq:moments}
  \E[\widehat\alpha^{\,2}]=\frac{N-1}{N}\alpha^{2},\qquad
  \Bias(\widehat\alpha^{\,2})=-\frac{\alpha^{2}}{N},\qquad
  \Var(\widehat\alpha^{\,2})=\frac{2(N-1)}{N^{2}}\alpha^{4},
\end{equation}
and consequently
\begin{equation}\label{eq:mse}
  \MSE(\widehat\alpha^{\,2})
  =\E\bigl[(\widehat\alpha^{\,2}-\alpha^{2})^{2}\bigr]=\frac{2N-1}{N^{2}}\alpha^{4}.
\end{equation}
Moreover $\E[\widehat\alpha]=b_N\alpha$ with
\begin{equation}\label{eq:bN}
  b_N=\sqrt{\frac2N}\,\frac{\Gamma(N/2)}{\Gamma((N-1)/2)}\in(0,1),\qquad b_N\uparrow1 .
\end{equation}
In particular $\widehat\alpha^{\,2}$ is biased but asymptotically unbiased, and
$\widetilde\alpha^{\,2}:=\frac{N}{N-1}\widehat\alpha^{\,2}$ is unbiased with
$\Var(\widetilde\alpha^{\,2})=2\alpha^{4}/(N-1)$.
\end{corollary}

\begin{proof}
If $Q\sim\chi^{2}_{N-1}$ then $\E Q=N-1$ and $\Var Q=2(N-1)$. By Theorem
\ref{th:exactlaw}(ii), $\widehat\alpha^{\,2}=\alpha^{2}Q/N$, which gives the three
identities in \eqref{eq:moments}. Then
\[
  \E\bigl[(\widehat\alpha^{\,2}-\alpha^{2})^{2}\bigr]
  =\Var(\widehat\alpha^{\,2})+\Bias(\widehat\alpha^{\,2})^{2}
  =\frac{2(N-1)}{N^{2}}\alpha^{4}+\frac{\alpha^{4}}{N^{2}}
  =\frac{2N-1}{N^{2}}\alpha^{4}.
\]
For \eqref{eq:bN}, $\widehat\alpha=\alpha\sqrt{Q/N}$ and
$\E\sqrt Q=\sqrt2\,\Gamma(N/2)/\Gamma((N-1)/2)$ for $Q\sim\chi^{2}_{N-1}$; that $b_N<1$
and $b_N\uparrow1$ follows from $\E\widehat\alpha\le(\E\widehat\alpha^{\,2})^{1/2}
=\alpha\sqrt{(N-1)/N}$ together with $b_N^{2}\to1$. The statements about
$\widetilde\alpha^{\,2}$ are immediate.
\end{proof}

\subsection{Optimality}

The exact laws of Theorem \ref{th:exactlaw} are accompanied by an exact optimality
theory, which we record next; it is a direct consequence of the exponential family
structure of \eqref{eq:Zlaw} and, through Remark \ref{re:glm}, of the classical theory of
the Gaussian linear model.

\begin{theorem}\label{th:umvu}
Under Assumption \ref{as:main} and for $N\ge2$:
\begin{enumerate}[label=\textup{(\roman*)}]
\item $\bigl(\widehat\theta,\widehat\alpha^{\,2}\bigr)$ is a complete sufficient
      statistic for $(\theta,\alpha^{2})\in\R\times(0,\infty)$;
\item $\widehat\theta$ is the uniformly minimum variance unbiased (UMVU) estimator of
      $\theta$;
\item $\widetilde\alpha^{\,2}=\frac{N}{N-1}\widehat\alpha^{\,2}$ is the UMVU estimator of
      $\alpha^{2}$, and $\widehat\alpha/b_N$, with $b_N$ as in \eqref{eq:bN}, is the UMVU
      estimator of $\alpha$.
\end{enumerate}
\end{theorem}

\begin{proof}
By \eqref{eq:loglik} the density of $\bZ$ can be written as
\[
  f(\bZ;\theta,\alpha^{2})=g(\theta,\alpha^{2})\,|\GH|^{-1/2}
  \exp\Bigl\{\eta_1\,\bZ^{\top}\GH^{-1}\bZ+\eta_2\,\bT^{\top}\GH^{-1}\bZ\Bigr\},
  \quad
  (\eta_1,\eta_2)=\Bigl(-\tfrac{1}{2\alpha^{2}},\tfrac{\theta}{\alpha^{2}}\Bigr),
\]
with $g(\theta,\alpha^{2})=(2\pi\alpha^{2})^{-N/2}
\exp\{-\theta^{2}d_N/(2\alpha^{2})\}$. This is a two-parameter exponential family with
sufficient statistic $S=(\bZ^{\top}\GH^{-1}\bZ,\ \bT^{\top}\GH^{-1}\bZ)$ and natural
parameter space $\{(\eta_1,\eta_2):\eta_1<0,\ \eta_2\in\R\}$, which contains a nonempty
open subset of $\R^{2}$. The family is therefore of full rank, and $S$ is complete and
sufficient \cite[Thm.~4.3.1 and Ch.~2]{LC98}. By \eqref{eq:mle} the map
$S\mapsto(\widehat\theta,\widehat\alpha^{\,2})
=\bigl(S_2/d_N,\ (S_1-S_2^{2}/d_N)/N\bigr)$ is a bijection of $\R^{2}$ onto its range,
so $(\widehat\theta,\widehat\alpha^{\,2})$ is complete and sufficient as well, proving
(i).

For (ii) and (iii), $\widehat\theta$, $\widetilde\alpha^{\,2}$ and $\widehat\alpha/b_N$
are functions of the complete sufficient statistic and are unbiased for $\theta$,
$\alpha^{2}$ and $\alpha$ respectively, by Theorem \ref{th:exactlaw}(i) and
Corollary \ref{co:moments}. The Lehmann--Scheff\'e theorem
\cite[Thm.~2.1.11]{LC98} then identifies each of them as the unique UMVU estimator of the
corresponding parameter.
\end{proof}

\begin{proposition}\label{pr:fisher}
Under Assumption \ref{as:main} the Fisher information matrix of
$\psi=(\theta,\alpha^{2})$ contained in $\bZ$ is
\begin{equation}\label{eq:fisher}
  I_N(\psi)=\begin{pmatrix} d_N/\alpha^{2} & 0\\[2pt] 0 & N/(2\alpha^{4})\end{pmatrix}.
\end{equation}
Consequently:
\begin{enumerate}[label=\textup{(\roman*)}]
\item $\theta$ and $\alpha^{2}$ are orthogonal parameters;
\item $\Var(\widehat\theta)=\alpha^{2}/d_N$, so $\widehat\theta$ attains the
      Cram\'er--Rao bound at every sample size $N\ge2$, not merely asymptotically;
\item the Cram\'er--Rao bound for unbiased estimation of $\alpha^{2}$ is
      $2\alpha^{4}/N$; it is attained by no unbiased estimator, and the UMVU estimator
      $\widetilde\alpha^{\,2}$ exceeds it by the factor $N/(N-1)\to1$.
\end{enumerate}
\end{proposition}

\begin{proof}
For $\bZ\sim\Norm(\mu(\psi),\Sigma(\psi))$ the information matrix has entries
$I_{ij}=(\partial_i\mu)^{\top}\Sigma^{-1}(\partial_j\mu)
+\tfrac12\tr\bigl(\Sigma^{-1}(\partial_i\Sigma)\Sigma^{-1}(\partial_j\Sigma)\bigr)$.
Here $\mu=\theta\bT$ and $\Sigma=\alpha^{2}\GH$, so $\partial_\theta\mu=\bT$,
$\partial_{\alpha^{2}}\mu=\bze$, $\partial_\theta\Sigma=0$ and
$\partial_{\alpha^{2}}\Sigma=\GH$. Hence
$I_{\theta\theta}=\bT^{\top}(\alpha^{2}\GH)^{-1}\bT=d_N/\alpha^{2}$,
$I_{\theta\alpha^{2}}=0$ and
$I_{\alpha^{2}\alpha^{2}}=\tfrac12\tr\bigl((\alpha^{-2}I_N)^{2}\bigr)=N/(2\alpha^{4})$,
which is \eqref{eq:fisher}; (i) and (ii) follow from \eqref{eq:fisher} and Theorem
\ref{th:exactlaw}(i). For (iii), the bound is $[I_N(\psi)^{-1}]_{22}=2\alpha^{4}/N$,
while $\Var(\widetilde\alpha^{\,2})=2\alpha^{4}/(N-1)$ by Corollary \ref{co:moments}.
Since $\widetilde\alpha^{\,2}$ has minimum variance among \emph{all} unbiased estimators
of $\alpha^{2}$ by Theorem \ref{th:umvu}(iii), and its variance is strictly larger than
$2\alpha^{4}/N$, no unbiased estimator attains the Cram\'er--Rao bound.
\end{proof}

\begin{remark}\label{re:msecomp}
Consider the one-parameter family $\{c\,\widehat\alpha^{\,2}:c>0\}$, which contains the
MLE ($c=1$) and the UMVU estimator ($c=N/(N-1)$). By Theorem \ref{th:exactlaw}(ii),
$\MSE(c\widehat\alpha^{\,2})
=\alpha^{4}\bigl[k^{2}(N-1)(N+1)-2k(N-1)+1\bigr]$ with $k=c/N$, which is minimised at
$k=1/(N+1)$, that is at $c=N/(N+1)$. Consequently
\begin{equation}\label{eq:msecomp}
  \MSE\Bigl(\tfrac{N}{N+1}\widehat\alpha^{\,2}\Bigr)=\frac{2\alpha^{4}}{N+1}
  <\MSE(\widehat\alpha^{\,2})=\frac{(2N-1)\alpha^{4}}{N^{2}}
  <\Var(\widetilde\alpha^{\,2})=\frac{2\alpha^{4}}{N-1}.
\end{equation}
Thus the maximum likelihood estimator dominates the UMVU estimator in mean square error,
and is itself dominated by $\frac{N}{N+1}\widehat\alpha^{\,2}$; the three estimators
agree to first order, all three mean square errors being $2\alpha^{4}/N+O(N^{-2})$. Which
one to prefer is therefore a matter of whether unbiasedness or mean square accuracy is
the design criterion.
\end{remark}

\subsection{Exact confidence intervals and exact tests}

\begin{corollary}\label{co:ci}
Let $\beta\in(0,1)$, let $\chi^{2}_{k,p}$ denote the $p$-quantile of the $\chi^{2}_{k}$
law and $t_{k,p}$ that of the Student law with $k$ degrees of freedom. Then, for every
$N\ge2$,
\begin{equation}\label{eq:ci-alpha}
  \Prob\left(\frac{N\widehat\alpha^{\,2}}{\chi^{2}_{N-1,\,1-\beta/2}}\le\alpha^{2}
  \le\frac{N\widehat\alpha^{\,2}}{\chi^{2}_{N-1,\,\beta/2}}\right)=1-\beta,
\end{equation}
and, with $\widetilde\alpha^{\,2}=\frac{N}{N-1}\widehat\alpha^{\,2}$,
\begin{equation}\label{eq:ci-theta}
  \sqrt{d_N}\;\frac{\widehat\theta-\theta}{\widetilde\alpha}\sim t_{N-1},
  \qquad\text{so that}\qquad
  \Prob\left(\bigl|\widehat\theta-\theta\bigr|
  \le t_{N-1,\,1-\beta/2}\,\frac{\widetilde\alpha}{\sqrt{d_N}}\right)=1-\beta .
\end{equation}
\end{corollary}

\begin{proof}
The interval for $\alpha^{2}$ is a direct inversion of Theorem \ref{th:exactlaw}(ii). For
the second statement, put $\xi=\sqrt{d_N}(\widehat\theta-\theta)/\alpha$ and
$Q=N\widehat\alpha^{\,2}/\alpha^{2}$. By Theorem \ref{th:exactlaw}, $\xi\sim\Norm(0,1)$
and $Q\sim\chi^{2}_{N-1}$ are independent, hence
$\xi/\sqrt{Q/(N-1)}\sim t_{N-1}$. Since $\sqrt{Q/(N-1)}=\widetilde\alpha/\alpha$, this
ratio equals $\sqrt{d_N}(\widehat\theta-\theta)/\widetilde\alpha$.
\end{proof}

\begin{corollary}\label{co:tests}
Fix $\beta\in(0,1)$ and $N\ge2$.
\begin{enumerate}[label=\textup{(\roman*)}]
\item For $\mathcal H_0:\theta=\theta_0$ against $\mathcal H_1:\theta\ne\theta_0$, the
      test that rejects when
      $\bigl|\sqrt{d_N}(\widehat\theta-\theta_0)/\widetilde\alpha\bigr|
      >t_{N-1,\,1-\beta/2}$ has exact size $\beta$, whatever $\alpha>0$, and is uniformly
      most powerful unbiased.
\item For $\mathcal H_0:\alpha^{2}=\alpha_0^{2}$ against
      $\mathcal H_1:\alpha^{2}\ne\alpha_0^{2}$, the test that rejects when
      $N\widehat\alpha^{\,2}/\alpha_0^{2}\notin
      [\chi^{2}_{N-1,\beta/2},\chi^{2}_{N-1,1-\beta/2}]$ has exact size $\beta$, whatever
      $\theta\in\R$.
\end{enumerate}
\end{corollary}

\begin{proof}
Exactness of the sizes is immediate from Corollary \ref{co:ci}. By Remark \ref{re:glm}
the model is, after an orthogonal change of coordinates, the Gaussian linear model with
one known regressor and unknown error variance, for which the equal-tailed $t$-test of a
linear hypothesis is uniformly most powerful unbiased
\cite[Ch.~5 and Sec.~7.1]{LR05}; this gives the last assertion of (i).
\end{proof}

\begin{remark}\label{re:cicomment}
Corollary \ref{co:ci} gives exact intervals at every sample size, with no appeal to
asymptotics, and the interval for $\theta$ is studentised, hence usable when $\alpha$ is
unknown. The only model-dependent quantity that must be computed is the scalar $d_N$,
obtained from a single Cholesky factorisation of $\GH$; by
Proposition \ref{pr:increment} it can equally be computed from the Toeplitz matrix
$C_N$, for which $O(N^{2})$ algorithms are available.
\end{remark}

\section{Asymptotic theory}\label{se:asymp}

Throughout this section the mesh $h>0$ is fixed and $N\to\infty$. We write
\[
\begin{gathered}
  \Hmax=\max_{1\le r\le m}H_r\in(1/2,1),\qquad
  \Hmin=\min_{1\le r\le m}H_r,\\[2pt]
  d_N=\bT_N^{\top}\Gamma_{H,N}^{-1}\bT_N,\qquad
  v_N=\Var\bigl(\widehat\theta^{(N)}\bigr)=\frac{\alpha^{2}}{d_N},
\end{gathered}
\]
and $\widehat\theta^{(N)},\widehat\alpha^{\,2}_N$ for the estimators based on
$\bZ^{(N)}=(Z(t_1),\dots,Z(t_N))^{\top}$. The letter $C$ denotes a finite positive
constant that may change from line to line and depends only on $h$, $m$, $H$, $\alpha$
and, where indicated, on a moment order $q$.

\subsection{A non-asymptotic variance bound for the drift estimator}

\begin{lemma}\label{le:kanto}
Let $A\in\R^{N\times N}$ be symmetric positive definite and $x\in\R^{N}$, $x\ne0$. Then
\begin{equation}\label{eq:kanto}
  x^{\top}A^{-1}x\ \ge\ \frac{\|x\|^{4}}{x^{\top}Ax}.
\end{equation}
\end{lemma}

\begin{proof}
Let $A^{1/2}$ be the symmetric positive definite square root of $A$ and
$A^{-1/2}=(A^{1/2})^{-1}$. Then
$\|x\|^{2}=x^{\top}x=\langle A^{1/2}x,A^{-1/2}x\rangle$, and Cauchy--Schwarz gives
\[
  \|x\|^{4}=\langle A^{1/2}x,A^{-1/2}x\rangle^{2}
  \le\|A^{1/2}x\|^{2}\|A^{-1/2}x\|^{2}=(x^{\top}Ax)(x^{\top}A^{-1}x).
\]
Since $x\ne0$ and $A$ is positive definite, $x^{\top}Ax>0$, and dividing by it yields
\eqref{eq:kanto}.
\end{proof}

\begin{remark}\label{re:kantorole}
Inequality \eqref{eq:kanto} is of the type used by Mishura, Ralchenko and Shklyar
\cite[Sec.~2]{MRS17}. Its role here is to convert a bound on the direct quadratic form
$\bT^{\top}\GH\bT$, which is explicitly computable from \eqref{eq:GammaHr}, into a lower
bound on the inverse form $d_N=\bT^{\top}\GH^{-1}\bT$, which is not.
\end{remark}

\begin{proposition}\label{pr:varbound}
For every $N\ge1$,
\begin{equation}\label{eq:varbound-sharp}
  \Var\bigl(\widehat\theta^{(N)}\bigr)\ \le\
  \frac{\alpha^{2}}{\|\bT\|^{4}}
  \Bigl(\sum_{j=1}^{N}t_j\Bigr)
  \Bigl(\sum_{r=1}^{m}\sum_{i=1}^{N}t_i^{2H_r+1}\Bigr)
\end{equation}
and consequently
\begin{equation}\label{eq:varbound}
  \Var\bigl(\widehat\theta^{(N)}\bigr)\ \le\
  \frac92\,\alpha^{2}\sum_{r=1}^{m}(Nh)^{2H_r-2}
  =\frac92\,\alpha^{2}\sum_{r=1}^{m}t_N^{2H_r-2}.
\end{equation}
If moreover $Nh\ge1$, then
$\Var(\widehat\theta^{(N)})\le\frac92m\alpha^{2}(Nh)^{2\Hmax-2}$. In particular
$\widehat\theta^{(N)}$ is unbiased with $\Var(\widehat\theta^{(N)})\to0$; that is,
$\widehat\theta^{(N)}\to\theta$ in mean square as $N\to\infty$.
\end{proposition}

\begin{proof}
Unbiasedness is Theorem \ref{th:exactlaw}(i). Combining
$\Var(\widehat\theta^{(N)})=\alpha^{2}/d_N$ with Lemma \ref{le:kanto} applied to
$A=\GH$ and $x=\bT$ gives
\begin{equation}\label{eq:varstep}
  \Var\bigl(\widehat\theta^{(N)}\bigr)=\frac{\alpha^{2}}{\bT^{\top}\GH^{-1}\bT}
  \le\frac{\alpha^{2}\,\bT^{\top}\GH\bT}{\|\bT\|^{4}} .
\end{equation}
We bound the numerator. Since $t_it_j\ge0$ and $|t_i-t_j|^{2H_r}\ge0$,
\eqref{eq:GammaHr} yields
\[
  \bT^{\top}\Gamma_{H_r}\bT=\sum_{i,j=1}^{N}t_it_j(\Gamma_{H_r})_{ij}
  \le\frac12\sum_{i,j=1}^{N}t_it_j\bigl(t_i^{2H_r}+t_j^{2H_r}\bigr)
  =\Bigl(\sum_{i=1}^{N}t_i^{2H_r+1}\Bigr)\Bigl(\sum_{j=1}^{N}t_j\Bigr),
\]
the last step by symmetry in $i$ and $j$. Summing over $r$ and inserting the result in
\eqref{eq:varstep} gives \eqref{eq:varbound-sharp}.

For \eqref{eq:varbound} we evaluate the elementary sums. With $t_k=kh$,
\[
  \sum_{j=1}^{N}t_j=\frac{hN(N+1)}{2},\qquad
  \|\bT\|^{2}=h^{2}\sum_{i=1}^{N}i^{2}=\frac{h^{2}N(N+1)(2N+1)}{6},
\]
so that, using $N+1\ge N$ and $(2N+1)^{2}\ge4N^{2}$,
\[
  \frac{\sum_{j=1}^{N}t_j}{\|\bT\|^{4}}
  =\frac{18}{h^{3}N(N+1)(2N+1)^{2}}\le\frac{18}{4h^{3}N^{4}}=\frac{9}{2h^{3}N^{4}} .
\]
Moreover $\sum_{i=1}^{N}t_i^{2H_r+1}=h^{2H_r+1}\sum_{i=1}^{N}i^{2H_r+1}
\le h^{2H_r+1}N^{2H_r+2}$. Substituting the last two displays into
\eqref{eq:varbound-sharp},
\[
  \Var\bigl(\widehat\theta^{(N)}\bigr)
  \le\frac{9\alpha^{2}}{2h^{3}N^{4}}\sum_{r=1}^{m}h^{2H_r+1}N^{2H_r+2}
  =\frac92\alpha^{2}\sum_{r=1}^{m}h^{2H_r-2}N^{2H_r-2},
\]
which is \eqref{eq:varbound}. If $Nh\ge1$, then since $2H_r-2\le2\Hmax-2$ we have
$(Nh)^{2H_r-2}\le(Nh)^{2\Hmax-2}$ for every $r$, giving the stated simplification.
Finally $2\Hmax-2<0$ forces $\Var(\widehat\theta^{(N)})\to0$, and combined with
unbiasedness this is mean square convergence.
\end{proof}

\begin{remark}\label{re:sharp}
By Proposition \ref{pr:increment}, $d_N=h^{2}\bone^{\top}C_N^{-1}\bone$, where $C_N$ is
the covariance matrix of the stationary Gaussian sequence $\bDelta$ of
Lemma \ref{le:statincr}. The spectral density of that sequence is
$f(\lambda)=\sum_{r=1}^{m}h^{2H_r}f_{H_r}(\lambda)$, $\lambda\in[-\pi,\pi]$, where
$f_{H}$ is the spectral density of standard fractional Gaussian noise and satisfies
$f_{H}(\lambda)\asymp|\lambda|^{1-2H}$ as $\lambda\to0$. Hence $f$ is regularly varying
at the origin with index $1-2\Hmax$. By the classical theory of generalised least
squares estimation of the mean of a long-memory stationary sequence, due to Adenstedt
\cite{Ade74}, one then has
$\bone^{\top}C_N^{-1}\bone\asymp h^{-2\Hmax}N^{2-2\Hmax}$, so that
\[
  \Var\bigl(\widehat\theta^{(N)}\bigr)\ \asymp\ (Nh)^{2\Hmax-2}.
\]
The bound \eqref{eq:varbound} therefore has the exact order in both $N$ and $h$, and only
its constant is conservative. A direct numerical evaluation of $\alpha^{2}/d_N$ for
$h\in\{1/252,1/12,1\}$ and $30\le N\le500$ with $H=(0.65,0.75,0.85)$ shows that the ratio
of the right-hand side of \eqref{eq:varbound} to $\Var(\widehat\theta^{(N)})$ stays
between $4.55$ and $4.57$ throughout, while the computable bound
\eqref{eq:varbound-sharp} stays within $31\%$ of the true variance. The approach to the
limiting exponent is, however, slow when the Hurst parameters are close: the components
with $H_r<\Hmax$ contribute a relative correction of order $N^{-2(\Hmax-H_r)}$, and for
$H=(0.65,0.75,0.85)$ the quantity $N^{2-2\Hmax}\Var(\widehat\theta^{(N)})$ still exceeds
its limiting value by about one third at $N=500$. This is why the effective exponent seen
at the sample sizes of Section \ref{se:sim} lies strictly between $2\Hmin-2$ and
$2\Hmax-2$.
\end{remark}

\begin{remark}\label{re:window}
Bound \eqref{eq:varbound} depends on $N$ and $h$ only through the length $t_N=Nh$ of the
observation window. This matches the classical picture for drift estimation: refining the
mesh on a fixed window does not, to this order, improve the drift estimate; only
extending the window does. It also explains the pattern seen in Section \ref{se:sim},
where for a fixed $N$ the dispersion of $\widehat\theta$ is markedly smaller under the
coarser mesh $h=1/12$ than under $h=1/252$. Finally, the component with the largest Hurst
parameter dominates the bound, so the strongest long-range dependence present in the
superposition governs the rate.
\end{remark}

\subsection{Strong consistency}
We use repeatedly the following form of Nelson's hypercontractivity inequality (see
\cite[Thm.~2.7.2]{NP12}): if $F$ belongs to the $p$-th Wiener chaos of an isonormal
Gaussian process, then for every $q\ge2$
\begin{equation}\label{eq:hyper}
  \E\bigl[|F|^{q}\bigr]\le(q-1)^{pq/2}\bigl(\E[F^{2}]\bigr)^{q/2}.
\end{equation}
Here all random variables are measurable with respect to the Gaussian space generated by
$\{B^{H_r}(t):t\ge0,\ r=1,\dots,m\}$, which we regard as an isonormal Gaussian process
over a separable Hilbert space $\Hi$.

\begin{theorem}\label{th:sctheta}
$\widehat\theta^{(N)}\to\theta$ almost surely as $N\to\infty$.
\end{theorem}

\begin{proof}
Fix $\gamma$ with $0<\gamma<1-\Hmax$, which is possible since $\Hmax<1$, and fix $q\ge2$
to be chosen below. By Theorem \ref{th:exactlaw}(i), $\widehat\theta^{(N)}-\theta$ is
centred Gaussian, hence lies in the first Wiener chaos and \eqref{eq:hyper} applies with
$p=1$ (equivalently, one may use the exact Gaussian moment formula). Together with
Proposition \ref{pr:varbound}, for $N\ge1/h$,
\[
  \E\bigl[|\widehat\theta^{(N)}-\theta|^{q}\bigr]
  \le(q-1)^{q/2}\bigl(\Var(\widehat\theta^{(N)})\bigr)^{q/2}
  \le C_qN^{q(\Hmax-1)} .
\]
By Markov's inequality,
\[
  \Prob\bigl(|\widehat\theta^{(N)}-\theta|>N^{-\gamma}\bigr)
  \le N^{q\gamma}\,\E\bigl[|\widehat\theta^{(N)}-\theta|^{q}\bigr]
  \le C_qN^{q(\gamma+\Hmax-1)} .
\]
Since $\gamma+\Hmax-1<0$, choosing $q$ so large that $q(\gamma+\Hmax-1)<-1$ makes the
right-hand side summable in $N$. By the Borel--Cantelli lemma,
$\Prob\bigl(|\widehat\theta^{(N)}-\theta|>N^{-\gamma}\ \text{i.o.}\bigr)=0$, so almost
surely $|\widehat\theta^{(N)}-\theta|\le N^{-\gamma}$ for all large $N$, and therefore
$\widehat\theta^{(N)}\to\theta$ almost surely.
\end{proof}

\begin{theorem}\label{th:scalpha}
$\widehat\alpha^{\,2}_N\to\alpha^{2}$ almost surely as $N\to\infty$.
\end{theorem}

\begin{proof}
By Lemma \ref{le:canon} and Theorem \ref{th:exactlaw}(ii), write
$Q_N=N\widehat\alpha^{\,2}_N/\alpha^{2}=\bV^{\top}P\bV\sim\chi^{2}_{N-1}$ and
\begin{equation}\label{eq:GN}
  \widehat\alpha^{\,2}_N-\alpha^{2}=\frac{\alpha^{2}}{N}\bigl(Q_N-N\bigr)
  =\frac{\alpha^{2}}{N}G_N-\frac{\alpha^{2}}{N},
  \qquad G_N:=Q_N-\tr P=Q_N-(N-1).
\end{equation}
Writing $V_i=W(e_i)$ for an orthonormal family $\{e_i\}_{i\le N}$ in $\Hi$, we have
$G_N=\sum_{i,j}P_{ij}(V_iV_j-\delta_{ij})
=I_2\bigl(\sum_{i,j}P_{ij}e_i\otimes e_j\bigr)$, an element of the second Wiener chaos,
with $\E[G_N^{2}]=\Var(Q_N)=2(N-1)$. By \eqref{eq:hyper} with $p=2$, for every $q\ge2$,
\[
  \E\bigl[|G_N|^{q}\bigr]\le(q-1)^{q}\bigl(2(N-1)\bigr)^{q/2}\le C_qN^{q/2}.
\]
Hence, by \eqref{eq:GN} and the elementary inequality
$|a+b|^{q}\le2^{q-1}(|a|^{q}+|b|^{q})$,
\[
  \E\bigl[|\widehat\alpha^{\,2}_N-\alpha^{2}|^{q}\bigr]
  \le\frac{2^{q-1}\alpha^{2q}}{N^{q}}\Bigl(\E\bigl[|G_N|^{q}\bigr]+1\Bigr)
  \le C_qN^{-q/2}.
\]
Fix $0<\delta<\frac12$. Markov's inequality gives
\[
  \Prob\bigl(|\widehat\alpha^{\,2}_N-\alpha^{2}|>N^{-\delta}\bigr)
  \le N^{q\delta}\,\E\bigl[|\widehat\alpha^{\,2}_N-\alpha^{2}|^{q}\bigr]
  \le C_qN^{q(\delta-\frac12)},
\]
which is summable in $N$ once $q$ is chosen with $q(\delta-\frac12)<-1$.
Borel--Cantelli then yields
$\Prob\bigl(|\widehat\alpha^{\,2}_N-\alpha^{2}|>N^{-\delta}\ \text{i.o.}\bigr)=0$ and
therefore $\widehat\alpha^{\,2}_N\to\alpha^{2}$ almost surely.
\end{proof}

\subsection{The structure of the sequence $\{\widehat\theta^{(N)}\}$}

\begin{theorem}\label{th:indincr}
The centred sequence $\xi_N:=\widehat\theta^{(N)}-\theta$, $N\ge1$, is jointly Gaussian
and satisfies
\begin{equation}\label{eq:covxi}
  \E\bigl[\xi_N\xi_{N'}\bigr]=v_{N\vee N'}\qquad\text{for all }N,N'\ge1 .
\end{equation}
Consequently $\{\widehat\theta^{(N)}\}_{N\ge1}$ has independent increments: for any
$N_1\le N_2\le N_3\le N_4$,
\[
  \Cov\bigl(\widehat\theta^{(N_4)}-\widehat\theta^{(N_3)},\
            \widehat\theta^{(N_2)}-\widehat\theta^{(N_1)}\bigr)=0,
\]
and the two increments are independent.
\end{theorem}

\begin{proof}
Each $\widehat\theta^{(N)}$ is a fixed linear functional of $\bZ^{(N)}$, hence of the
underlying Gaussian family; therefore $\{\xi_N\}_{N\ge1}$ is jointly Gaussian and it
suffices to compute covariances.

Fix $N_2\le N_3$ and let $J=\bigl(I_{N_2},\,0_{N_2\times(N_3-N_2)}\bigr)
\in\R^{N_2\times N_3}$, so that
\[
  \bZ^{(N_2)}=J\bZ^{(N_3)},\qquad \bT_{N_2}=J\bT_{N_3},\qquad
  \Gamma_{H,N_2}=J\Gamma_{H,N_3}J^{\top} .
\]
Writing $\bW^{(N)}=\bZ^{(N)}-\theta\bT_N$, we have
\[
  \Cov\bigl(\bW^{(N_3)},\bW^{(N_2)}\bigr)=\alpha^{2}\Gamma_{H,N_3}J^{\top},
  \qquad
  \xi_N=\frac{\bT_N^{\top}\Gamma_{H,N}^{-1}\bW^{(N)}}{d_N}.
\]
Hence
\[
  \E[\xi_{N_3}\xi_{N_2}]
  =\frac{\bT_{N_3}^{\top}\Gamma_{H,N_3}^{-1}
        \bigl(\alpha^{2}\Gamma_{H,N_3}J^{\top}\bigr)\Gamma_{H,N_2}^{-1}\bT_{N_2}}
        {d_{N_3}d_{N_2}}
  =\frac{\alpha^{2}(J\bT_{N_3})^{\top}\Gamma_{H,N_2}^{-1}\bT_{N_2}}{d_{N_3}d_{N_2}} .
\]
Since $J\bT_{N_3}=\bT_{N_2}$, the numerator equals $\alpha^{2}d_{N_2}$, so
$\E[\xi_{N_3}\xi_{N_2}]=\alpha^{2}/d_{N_3}=v_{N_3}$, which is \eqref{eq:covxi}.

Now let $N_1\le N_2\le N_3\le N_4$. Expanding and using \eqref{eq:covxi},
\[
\begin{aligned}
  \Cov\bigl(\xi_{N_4}-\xi_{N_3},\ \xi_{N_2}-\xi_{N_1}\bigr)
  &=\E[\xi_{N_4}\xi_{N_2}]-\E[\xi_{N_4}\xi_{N_1}]
    -\E[\xi_{N_3}\xi_{N_2}]+\E[\xi_{N_3}\xi_{N_1}]\\
  &=v_{N_4}-v_{N_4}-v_{N_3}+v_{N_3}=0 .
\end{aligned}
\]
The increments $\xi_{N_4}-\xi_{N_3}$ and $\xi_{N_2}-\xi_{N_1}$ are jointly Gaussian and
uncorrelated, hence independent; the same computation applied to finitely many pairwise
disjoint index blocks gives pairwise uncorrelatedness of the corresponding increments and
therefore, by joint Gaussianity, their mutual independence.
\end{proof}

\begin{corollary}\label{co:tcbm}
The variance sequence $(v_N)_{N\ge1}$ is nonincreasing with $v_N\downarrow0$, and there
exists a standard Brownian motion $\{W(v),\,v\ge0\}$ with $W(0)=0$ such that
\[
  \bigl(\widehat\theta^{(N)}-\theta\bigr)_{N\ge1}
  \overset{d}{=}\bigl(W(v_N)\bigr)_{N\ge1}
\]
as random elements of the sequence space $\R^{\N}$ equipped with its product
$\sigma$-field.
\end{corollary}

\begin{proof}
Expanding $\Var(\xi_{N+1}-\xi_N)$ and using \eqref{eq:covxi}, which gives
$\E[\xi_N\xi_{N+1}]=v_{N+1}$,
\[
  0\le\Var(\xi_{N+1}-\xi_N)=v_{N+1}-2v_{N+1}+v_N=v_N-v_{N+1},
\]
so $v_{N+1}\le v_N$. By Proposition \ref{pr:varbound}, $v_N\to0$, hence
$v_N\downarrow0$.

Let $W$ be a standard Brownian motion. Both $(\xi_N)_{N\ge1}$ and
$(W(v_N))_{N\ge1}$ are centred Gaussian sequences, and since $(v_N)$ is nonincreasing,
$\Cov\bigl(W(v_N),W(v_{N'})\bigr)=v_N\wedge v_{N'}=v_{N\vee N'}$, which by
\eqref{eq:covxi} is exactly $\E[\xi_N\xi_{N'}]$. Two centred Gaussian sequences with the
same covariance have the same law on $\R^{\N}$.
\end{proof}

\begin{remark}\label{re:secondproof}
Corollary \ref{co:tcbm} yields strong consistency of $\widehat\theta$ without any moment
computation. Indeed, the set $\Lambda=\{x\in\R^{\N}:\lim_Nx_N=0\}$ belongs to the product
$\sigma$-field of $\R^{\N}$, so
$\Prob(\lim_N\xi_N=0)=\Prob(\lim_NW(v_N)=0)$. Since $v_N\downarrow0$ and $W$ is almost
surely continuous at $0$ with $W(0)=0$, the right-hand side equals $1$.
\end{remark}

\begin{remark}\label{re:indincrcomment}
Theorem \ref{th:indincr} is noteworthy because the underlying data are strongly
dependent: the increments of the estimator sequence are independent even though the
increments of the observed process are not. The mechanism is the projection structure of
the generalised least squares estimator combined with the nestedness of the observation
vectors, which makes $\widehat\theta^{(N)}-\theta$ behave exactly like a Brownian motion
sampled along its own (decreasing) variance scale.
\end{remark}

\subsection{Asymptotic normality}

\begin{theorem}\label{th:an}
As $N\to\infty$,
\begin{equation}\label{eq:an-theta}
  \sqrt{d_N}\,\bigl(\widehat\theta^{(N)}-\theta\bigr)\sim\Norm(0,\alpha^{2})
  \quad\text{for every }N,\qquad\text{hence}\qquad
  \sqrt{d_N}\,\bigl(\widehat\theta^{(N)}-\theta\bigr)\xrightarrow{\ d\ }
  \Norm(0,\alpha^{2}),
\end{equation}
and
\begin{equation}\label{eq:an-alpha}
  F_N:=\frac{1}{\alpha^{2}}\sqrt{\frac N2}
  \bigl(\widehat\alpha^{\,2}_N-\alpha^{2}\bigr)\xrightarrow{\ d\ }\Norm(0,1).
\end{equation}
Moreover, by Theorem \ref{th:exactlaw}(iii) the two statistics in \eqref{eq:an-theta} and
\eqref{eq:an-alpha} are independent for every $N$, so the convergence holds jointly with
a diagonal limiting covariance.
\end{theorem}

\begin{proof}
Statement \eqref{eq:an-theta} is immediate from Theorem \ref{th:exactlaw}(i): the
normalised drift estimator is exactly $\Norm(0,\alpha^{2})$ at every sample size. (The
normalisation is nondegenerate since $d_N=\alpha^{2}/v_N\to\infty$ by
Proposition \ref{pr:varbound}.)

For \eqref{eq:an-alpha}, use \eqref{eq:GN}: with $Q_N\sim\chi^{2}_{N-1}$,
\[
  F_N=\frac{1}{\sqrt{2N}}\bigl(Q_N-N\bigr)
  =\sqrt{\frac{N-1}{N}}\cdot\frac{Q_N-(N-1)}{\sqrt{2(N-1)}}-\frac{1}{\sqrt{2N}} .
\]
Since $Q_N$ is a sum of $N-1$ independent $\chi^{2}_{1}$ variables with mean $1$ and
variance $2$, the classical central limit theorem gives
$\bigl(Q_N-(N-1)\bigr)/\sqrt{2(N-1)}\xrightarrow{d}\Norm(0,1)$. As $(N-1)/N\to1$ and
$(2N)^{-1/2}\to0$, Slutsky's theorem yields \eqref{eq:an-alpha}.

Finally, joint convergence with a diagonal limit follows from the exact independence of
$\widehat\theta^{(N)}$ and $\widehat\alpha^{\,2}_N$ established in Theorem
\ref{th:exactlaw}(iii).
\end{proof}

The elementary proof above exploits the exact chi-square law of Theorem
\ref{th:exactlaw}, and in that sense the second proof below is not needed here. We record
it because it is the argument that survives the relaxation of the common-scale
assumption: with component-specific scales the exact chi-square law fails, while the
normalised scale estimator remains in the second Wiener chaos and the criterion of
Nualart and Ortiz-Latorre \cite{NOL08} still applies. We use the
notation of the proof of Theorem \ref{th:scalpha}.

\begin{lemma}\label{le:malliavin}
Let $G_N=N\widehat\alpha^{\,2}_N/\alpha^{2}-(N-1)$ and
$\widetilde G_N=G_N/\sqrt{2N}$, so that $F_N=\widetilde G_N-(2N)^{-1/2}$. Then
$\widetilde G_N$ belongs to the second Wiener chaos,
$\E[\widetilde G_N^{2}]=(N-1)/N$, and
\begin{equation}\label{eq:mall}
  \bigl\|D\widetilde G_N\bigr\|_{\Hi}^{2}=\frac{2\widehat\alpha^{\,2}_N}{\alpha^{2}} .
\end{equation}
\end{lemma}

\begin{proof}
By Lemma \ref{le:canon}, $N\widehat\alpha^{\,2}_N/\alpha^{2}=\bV^{\top}P\bV$ with
$\bV\sim\Norm(\bze,I_N)$ and $P$ an orthogonal projection of rank $N-1$. Writing
$V_i=W(e_i)$ for an orthonormal family $\{e_i\}_{i\le N}\subset\Hi$,
\[
  \widetilde G_N=\frac{1}{\sqrt{2N}}\sum_{i,j=1}^{N}P_{ij}\bigl(V_iV_j-\delta_{ij}\bigr)
  =\frac{1}{\sqrt{2N}}I_2\Bigl(\sum_{i,j=1}^{N}P_{ij}\,e_i\otimes e_j\Bigr),
\]
an element of the second chaos, with
$\E[\widetilde G_N^{2}]=\Var(\bV^{\top}P\bV)/(2N)=2\tr(P^{2})/(2N)=(N-1)/N$.
Differentiating,
\[
  D\widetilde G_N=\frac{1}{\sqrt{2N}}\sum_{i,j=1}^{N}P_{ij}\bigl(V_ie_j+V_je_i\bigr)
  =\frac{2}{\sqrt{2N}}\sum_{j=1}^{N}(P\bV)_je_j,
\]
using $P=P^{\top}$. Since the $e_j$ are orthonormal,
\[
  \|D\widetilde G_N\|^{2}_{\Hi}=\frac{4}{2N}\|P\bV\|^{2}
  =\frac2N\bV^{\top}P\bV=\frac{2\widehat\alpha^{\,2}_N}{\alpha^{2}} . \qedhere
\]
\end{proof}

\begin{proof}[Second proof of \eqref{eq:an-alpha}]
By Lemma \ref{le:malliavin}, $\widetilde G_N$ lies in the second Wiener chaos with
$\E[\widetilde G_N^{2}]\to1$. Moreover, by Corollary \ref{co:moments},
$\E[(\widehat\alpha^{\,2}_N-\alpha^{2})^{2}]=(2N-1)\alpha^{4}/N^{2}\to0$, that is
$\widehat\alpha^{\,2}_N\to\alpha^{2}$ in $L^{2}$, so \eqref{eq:mall} gives
\[
  \bigl\|D\widetilde G_N\bigr\|^{2}_{\Hi}
  =\frac{2\widehat\alpha^{\,2}_N}{\alpha^{2}}\longrightarrow2
  \qquad\text{in }L^{2}(\Omega).
\]
By the criterion of Nualart and Ortiz-Latorre \cite[Thm.~4]{NOL08} for a sequence in the
$p$-th chaos with $p=2$ and limiting variance $\sigma^{2}=1$ (convergence of
$\|D\widetilde G_N\|^{2}_{\Hi}$ in $L^{2}$ to $p\sigma^{2}=2$), we conclude
$\widetilde G_N\xrightarrow{d}\Norm(0,1)$, and hence
$F_N=\widetilde G_N-(2N)^{-1/2}\xrightarrow{d}\Norm(0,1)$.
\end{proof}

\begin{remark}\label{re:special}
By Remark \ref{re:half}, the results of Sections \ref{se:mle} and \ref{se:asymp}
remain valid for arbitrary $H_r\in(0,1)$. In particular the framework contains the
following models as special cases:
\begin{enumerate}[label=\textup{(\alph*)}]
\item $m=1$, $H_1=\frac12$: Brownian motion with drift;
\item $m=1$, $H_1\in(\frac12,1)$: the drifted fractional Brownian motion model of
      Hu et al.\ \cite{HNXZ11};
\item $m=2$, $H_1=\frac12$, $H_2\in(\frac12,1)$: the mixed fractional Brownian motion
      model of Xiao et al.\ \cite{XZZ11};
\item $m=2$, $H_1,H_2\in(\frac12,1)$, $H_1\ne H_2$: the double-fractional model of
      Mishura and Voronov \cite{MV15}.
\end{enumerate}
\end{remark}

\section{Simulation study}\label{se:sim}

\subsection{Design}
Sample paths are generated directly from the covariance matrix $\alpha^{2}\GH$ of the
model by Cholesky factorisation, which is exact: if $\GH=LL^{\top}$ with $L$ lower
triangular, then $\bZ=\theta\bT+\alpha L\varepsilon$ with
$\varepsilon\sim\Norm(\bze,I_N)$ has the exact law \eqref{eq:Zlaw}; see, e.g.,
\cite[Ch.~XI]{AG07} and \cite{Hig01}. The same factorisation is reused to evaluate
$\GH^{-1}$ in \eqref{eq:mle}, so that the estimators are computed without forming the
inverse explicitly. Throughout,
\[
  H_1=0.65,\qquad H_2=0.75,\qquad H_3=0.85\qquad(m=3),
\]
the sample size $N$ takes the values $\{30,100,200,300,500\}$, and the two sampling
frequencies $h=1/252$ (daily) and $h=1/12$ (monthly) are considered. Three parameter
scenarios are used, $(\theta,\alpha)\in\{(0.5,0.4),(1.5,1),(3,2)\}$.

Section \ref{sse:point} assesses point estimation and uses $500$ replications per
configuration, as is customary for the estimation of means and mean squared errors.
Section \ref{sse:exact} assesses the exact distributional theory (coverage
probabilities, goodness of fit) and uses $20\,000$ replications per configuration,
since discriminating between an exact and an approximately exact sampling distribution
requires substantially greater Monte Carlo precision. With $20\,000$ replications the
standard error of an estimated $0.95$ coverage probability is $0.0015$.

\begin{figure}[htbp]
\centering
\includegraphics[width=\textwidth]{fig1.pdf}
\caption{The multi-mixed fractional Brownian motion model \eqref{eq:model} for
Scenario 2 ($\theta=1.5$, $\alpha=1$) with $N=500$ and $h=1/252$, so that the observation
window is $t_N\approx1.98$. (a) Six independent trajectories and the trend $\theta t$.
(b) The three fractional components $\alpha B^{H_r}$ of one trajectory.}
\label{fi:paths}
\end{figure}

\subsection{Point estimation}\label{sse:point}
Tables \ref{ta:daily} and \ref{ta:monthly} report, for each configuration, the Monte
Carlo mean, bias, standard deviation and mean squared error, together with the
\emph{exact} standard deviation predicted by the theory. For ease of comparison with the
true value $\alpha$, the tables report the scale estimator on the natural scale,
$\widehat\alpha=\sqrt{\widehat\alpha^{\,2}}$; the exact entries are then
$\alpha/\sqrt{d_N}$ for $\widehat\theta$, and
$\bigl(\alpha^{2}(N-1)/N-b_N^{2}\alpha^{2}\bigr)^{1/2}$ for $\widehat\alpha$, with $b_N$
as in \eqref{eq:bN}.

The empirical means of $\widehat\theta$ and $\widehat\alpha$ remain close to the true
values in every scenario, and the standard deviations and mean squared errors decrease
monotonically in $N$, as predicted by Proposition \ref{pr:varbound} and
Corollary \ref{co:moments}. Beyond that, four features deserve comment, each of them a
direct check on the theory.

First, every simulated standard deviation agrees with the exact value to within Monte
Carlo error; with $500$ replications the relative standard error of an estimated standard
deviation is $(2\times499)^{-1/2}=3.2\%$, and across the sixty entries of
Tables \ref{ta:daily} and \ref{ta:monthly} the largest discrepancy is $4.8\%$, that is
$1.5$ Monte Carlo standard errors. The exact finite-sample theory is therefore not merely
qualitatively but quantitatively correct already at $N=30$.

Second, the scale estimator is accurate even at $N=30$, and its dispersion is
\emph{identical} under the two sampling frequencies: the exact standard deviation of
$\widehat\alpha$ is the same number in Table \ref{ta:daily} and in
Table \ref{ta:monthly}, row by row, and the simulated values agree to within Monte Carlo
error; for instance, in Scenario 2 with $N=500$ the simulated standard deviation is
$0.0323$ under both meshes, against the exact value $0.0316$. This is exactly what
Theorem \ref{th:exactlaw}(ii) predicts, since the law of
$\widehat\alpha^{\,2}/\alpha^{2}$ does not depend on $h$ (Remark \ref{re:pivot}).
Likewise the exact mean $\E[\widehat\alpha]=b_N\alpha$ equals $0.3899$ for $\alpha=0.4$,
$N=30$, against the simulated $0.3888$.

Third, the drift estimator is markedly more variable, and its dispersion scales linearly
in $\alpha$, as \eqref{eq:varbound} requires: in Table \ref{ta:daily} at $N=500$ the
standard deviations $0.5920$, $1.4470$ and $2.9366$ across the three scenarios are in the
ratios $1:2.44:4.96$, against the theoretical $\alpha$-ratios $1:2.5:5$.

Fourth, the comparison of the two sampling frequencies isolates the role of the
observation window. Under $h=1/12$ the standard deviation of $\widehat\theta$ is close to
half its value under $h=1/252$ at every sample size, whereas $\widehat\alpha$ is
unaffected. This is the content of Remark \ref{re:window}: the accuracy of drift
estimation is governed by the window length $t_N=Nh$, which is $21$ times longer under
monthly sampling. The exact ratio $\sqrt{d_N(1/252)/d_N(1/12)}$ implied by the theory
rises from $0.456$ at $N=30$ to $0.509$ at $N=500$, in close agreement with the simulated
ratios, which range from $0.44$ to $0.52$. It is worth noting that this ratio has not yet
reached its asymptotic value $21^{\Hmax-1}=0.633$: as explained in Remark \ref{re:sharp},
the sub-dominant components decay only polynomially and the effective exponent at these
sample sizes is still strictly between $\Hmin-1$ and $\Hmax-1$.

\begin{table}[htbp]
\caption{Monte Carlo results for the multi-mixed fBm model, $h=1/252$; $500$
replications, $m=3$, $(H_1,H_2,H_3)=(0.65,0.75,0.85)$. ``Exact s.d.'' is the standard
deviation predicted by Theorem \ref{th:exactlaw} and Corollary \ref{co:moments}.}
\label{ta:daily}
\small
\begin{tabular}{@{}llrrrrr@{}}
\toprule
$N$ & Parameter & Mean & Bias & Std.\ dev. & Exact s.d. & MSE\\
\midrule
\multicolumn{7}{@{}l}{\itshape Scenario 1: $\theta=0.5$, $\alpha=0.4$}\\
30  & $\widehat\theta$  & 0.5420 & \phantom{$-$}0.0420 & 1.1631 & 1.2086 & 1.3533\\
    & $\widehat\alpha$  & 0.3888 & $-$0.0112 & 0.0511 & 0.0514 & 0.0027\\
100 & $\widehat\theta$  & 0.4843 & $-$0.0157 & 0.8563 & 0.8727 & 0.7328\\
    & $\widehat\alpha$  & 0.3965 & $-$0.0035 & 0.0285 & 0.0282 & 0.0008\\
200 & $\widehat\theta$  & 0.5030 & \phantom{$-$}0.0030 & 0.7114 & 0.7298 & 0.5056\\
    & $\widehat\alpha$  & 0.3989 & $-$0.0011 & 0.0198 & 0.0200 & 0.0004\\
300 & $\widehat\theta$  & 0.5163 & \phantom{$-$}0.0163 & 0.6464 & 0.6593 & 0.4176\\
    & $\widehat\alpha$  & 0.3995 & $-$0.0005 & 0.0161 & 0.0163 & 0.0003\\
500 & $\widehat\theta$  & 0.4991 & $-$0.0009 & 0.5920 & 0.5819 & 0.3501\\
    & $\widehat\alpha$  & 0.3994 & $-$0.0006 & 0.0130 & 0.0126 & 0.0002\\
\addlinespace
\multicolumn{7}{@{}l}{\itshape Scenario 2: $\theta=1.5$, $\alpha=1.0$}\\
30  & $\widehat\theta$  & 1.5614 & \phantom{$-$}0.0614 & 3.1215 & 3.0215 & 9.7381\\
    & $\widehat\alpha$  & 0.9728 & $-$0.0272 & 0.1288 & 0.1285 & 0.0173\\
100 & $\widehat\theta$  & 1.5406 & \phantom{$-$}0.0406 & 2.0798 & 2.1818 & 4.3229\\
    & $\widehat\alpha$  & 0.9945 & $-$0.0055 & 0.0698 & 0.0706 & 0.0049\\
200 & $\widehat\theta$  & 1.3670 & $-$0.1330 & 1.7955 & 1.8245 & 3.2384\\
    & $\widehat\alpha$  & 0.9983 & $-$0.0017 & 0.0524 & 0.0500 & 0.0027\\
300 & $\widehat\theta$  & 1.4937 & $-$0.0063 & 1.6793 & 1.6482 & 2.8171\\
    & $\widehat\alpha$  & 0.9968 & $-$0.0032 & 0.0416 & 0.0408 & 0.0017\\
500 & $\widehat\theta$  & 1.4834 & $-$0.0166 & 1.4470 & 1.4547 & 2.0919\\
    & $\widehat\alpha$  & 0.9994 & $-$0.0006 & 0.0323 & 0.0316 & 0.0010\\
\addlinespace
\multicolumn{7}{@{}l}{\itshape Scenario 3: $\theta=3.0$, $\alpha=2.0$}\\
30  & $\widehat\theta$  & 3.1144 & \phantom{$-$}0.1144 & 5.9432 & 6.0431 & 35.3000\\
    & $\widehat\alpha$  & 1.9587 & $-$0.0413 & 0.2651 & 0.2571 & 0.0719\\
100 & $\widehat\theta$  & 3.2531 & \phantom{$-$}0.2531 & 4.4664 & 4.3636 & 19.9930\\
    & $\widehat\alpha$  & 1.9884 & $-$0.0116 & 0.1414 & 0.1412 & 0.0201\\
200 & $\widehat\theta$  & 3.0443 & \phantom{$-$}0.0443 & 3.6612 & 3.6491 & 13.3930\\
    & $\widehat\alpha$  & 1.9927 & $-$0.0073 & 0.1001 & 0.0999 & 0.0101\\
300 & $\widehat\theta$  & 2.9960 & $-$0.0040 & 3.2454 & 3.2964 & 10.5220\\
    & $\widehat\alpha$  & 1.9987 & $-$0.0013 & 0.0843 & 0.0816 & 0.0071\\
500 & $\widehat\theta$  & 2.9867 & $-$0.0133 & 2.9366 & 2.9093 & 8.6152\\
    & $\widehat\alpha$  & 1.9990 & $-$0.0010 & 0.0627 & 0.0632 & 0.0039\\
\bottomrule
\end{tabular}
\end{table}

\begin{table}[htbp]
\caption{Monte Carlo results for the multi-mixed fBm model, $h=1/12$; $500$
replications, $m=3$, $(H_1,H_2,H_3)=(0.65,0.75,0.85)$. The ``Exact s.d.'' entries for
$\widehat\alpha$ coincide with those of Table \ref{ta:daily}, as they must by
Theorem \ref{th:exactlaw}(ii).}
\label{ta:monthly}
\small
\begin{tabular}{@{}llrrrrr@{}}
\toprule
$N$ & Parameter & Mean & Bias & Std.\ dev. & Exact s.d. & MSE\\
\midrule
\multicolumn{7}{@{}l}{\itshape Scenario 1: $\theta=0.5$, $\alpha=0.4$}\\
30  & $\widehat\theta$  & 0.5038 & \phantom{$-$}0.0038 & 0.5610 & 0.5507 & 0.3144\\
    & $\widehat\alpha$  & 0.3901 & $-$0.0099 & 0.0520 & 0.0514 & 0.0028\\
100 & $\widehat\theta$  & 0.5120 & \phantom{$-$}0.0120 & 0.4178 & 0.4172 & 0.1745\\
    & $\widehat\alpha$  & 0.3955 & $-$0.0045 & 0.0284 & 0.0282 & 0.0008\\
200 & $\widehat\theta$  & 0.4990 & $-$0.0010 & 0.3656 & 0.3585 & 0.1336\\
    & $\widehat\alpha$  & 0.3982 & $-$0.0018 & 0.0207 & 0.0200 & 0.0004\\
300 & $\widehat\theta$  & 0.5154 & \phantom{$-$}0.0154 & 0.3264 & 0.3289 & 0.1067\\
    & $\widehat\alpha$  & 0.3986 & $-$0.0014 & 0.0160 & 0.0163 & 0.0003\\
500 & $\widehat\theta$  & 0.5137 & \phantom{$-$}0.0137 & 0.2904 & 0.2959 & 0.0844\\
    & $\widehat\alpha$  & 0.3999 & $-$0.0001 & 0.0125 & 0.0126 & 0.0002\\
\addlinespace
\multicolumn{7}{@{}l}{\itshape Scenario 2: $\theta=1.5$, $\alpha=1.0$}\\
30  & $\widehat\theta$  & 1.5184 & \phantom{$-$}0.0184 & 1.3724 & 1.3769 & 1.8820\\
    & $\widehat\alpha$  & 0.9778 & $-$0.0222 & 0.1328 & 0.1285 & 0.0181\\
100 & $\widehat\theta$  & 1.5246 & \phantom{$-$}0.0246 & 1.0378 & 1.0430 & 1.0765\\
    & $\widehat\alpha$  & 0.9916 & $-$0.0084 & 0.0725 & 0.0706 & 0.0053\\
200 & $\widehat\theta$  & 1.4817 & $-$0.0183 & 0.9138 & 0.8963 & 0.8345\\
    & $\widehat\alpha$  & 0.9969 & $-$0.0031 & 0.0505 & 0.0500 & 0.0026\\
300 & $\widehat\theta$  & 1.4945 & $-$0.0055 & 0.8237 & 0.8223 & 0.6778\\
    & $\widehat\alpha$  & 0.9978 & $-$0.0022 & 0.0396 & 0.0408 & 0.0016\\
500 & $\widehat\theta$  & 1.5111 & \phantom{$-$}0.0111 & 0.7477 & 0.7397 & 0.5585\\
    & $\widehat\alpha$  & 0.9984 & $-$0.0016 & 0.0323 & 0.0316 & 0.0010\\
\addlinespace
\multicolumn{7}{@{}l}{\itshape Scenario 3: $\theta=3.0$, $\alpha=2.0$}\\
30  & $\widehat\theta$  & 3.1339 & \phantom{$-$}0.1339 & 2.7766 & 2.7537 & 7.7196\\
    & $\widehat\alpha$  & 1.9432 & $-$0.0568 & 0.2588 & 0.2571 & 0.0701\\
100 & $\widehat\theta$  & 3.0232 & \phantom{$-$}0.0232 & 2.1174 & 2.0861 & 4.4793\\
    & $\widehat\alpha$  & 1.9872 & $-$0.0128 & 0.1405 & 0.1412 & 0.0199\\
200 & $\widehat\theta$  & 2.9658 & $-$0.0342 & 1.7160 & 1.7926 & 2.9430\\
    & $\widehat\alpha$  & 1.9955 & $-$0.0045 & 0.0993 & 0.0999 & 0.0099\\
300 & $\widehat\theta$  & 3.0484 & \phantom{$-$}0.0484 & 1.6349 & 1.6447 & 2.6724\\
    & $\widehat\alpha$  & 1.9933 & $-$0.0067 & 0.0821 & 0.0816 & 0.0068\\
500 & $\widehat\theta$  & 3.0122 & \phantom{$-$}0.0122 & 1.4436 & 1.4793 & 2.0820\\
    & $\widehat\alpha$  & 1.9990 & $-$0.0010 & 0.0612 & 0.0632 & 0.0037\\
\bottomrule
\end{tabular}
\end{table}

\subsection{Exactness of the finite-sample theory}\label{sse:exact}

Theorem \ref{th:exactlaw} and Corollary \ref{co:ci} make claims of a different nature
from those assessed above: they assert that the confidence intervals
\eqref{eq:ci-alpha} and \eqref{eq:ci-theta} have coverage exactly $1-\beta$ at every
sample size, and that the two estimators are exactly independent. Table \ref{ta:coverage}
reports the corresponding Monte Carlo check, based on $20\,000$ replications of
Scenario 2. By Remark \ref{re:pivot} the coverage of \eqref{eq:ci-alpha} does not depend
on $\theta$, $\alpha$, $h$ or $H$ at all, and the coverage of \eqref{eq:ci-theta} does
not depend on them either, so a single scenario suffices; the two meshes are retained as
a control.

Every entry of the two coverage columns lies within three Monte Carlo standard errors of
the nominal $0.95$, with no systematic drift in $N$ and, in particular, no deterioration
at $N=30$. The Kolmogorov--Smirnov tests of the two pivotal statistics against their
exact reference laws ($t_{N-1}$ for
$\sqrt{d_N}(\widehat\theta-\theta)/\widetilde\alpha$ and $\chi^{2}_{N-1}$ for
$N\widehat\alpha^{\,2}/\alpha^{2}$) produce no evidence against the theory: of the
twenty $p$-values reported, one falls below $0.05$, which is what one expects by chance.
The empirical correlation between $\widehat\theta$ and $\widehat\alpha^{\,2}$ never
exceeds $0.010$ in absolute value, against a Monte Carlo standard error of
$1/\sqrt{20\,000}=0.007$, confirming Theorem \ref{th:exactlaw}(iii).

\begin{figure}[htbp]
\centering
\includegraphics[width=\textwidth]{fig2.pdf}
\caption{Histograms of $\widehat\theta$ (top) and $\widehat\alpha^{\,2}$ (bottom) over
$20\,000$ replications of Scenario 2 with $h=1/252$, at $N=30$ (left) and $N=500$
(right), with the exact densities of Theorem \ref{th:exactlaw} superimposed: the
$\Norm(\theta,\alpha^{2}/d_N)$ density and the density of
$(\alpha^{2}/N)\chi^{2}_{N-1}$. Dotted vertical lines mark the true parameter values.}
\label{fi:hist}
\end{figure}

\begin{figure}[htbp]
\centering
\includegraphics[width=\textwidth]{fig3.pdf}
\caption{Quantile--quantile plots against the exact reference laws of
Theorem \ref{th:exactlaw}, for the same $20\,000$ replications as in
Figure \ref{fi:hist}: $\sqrt{d_N}(\widehat\theta-\theta)/\alpha$ against $\Norm(0,1)$
(top) and $N\widehat\alpha^{\,2}/\alpha^{2}$ against $\chi^{2}_{N-1}$ (bottom), at
$N=30$ (left) and $N=500$ (right).}
\label{fi:qq}
\end{figure}

\begin{table}[htbp]
\caption{Exactness of the finite-sample theory: empirical coverage of the nominal $95\%$
intervals \eqref{eq:ci-theta} and \eqref{eq:ci-alpha}, Kolmogorov--Smirnov $p$-values of
the two pivotal statistics against their exact reference laws, and empirical correlation
$\widehat\rho$ between $\widehat\theta$ and $\widehat\alpha^{\,2}$. Scenario 2
($\theta=1.5$, $\alpha=1$), $m=3$,
$(H_1,H_2,H_3)=(0.65,0.75,0.85)$, $20\,000$ replications; the Monte Carlo standard error
of a coverage entry is $0.0015$.}
\label{ta:coverage}
\small
\begin{tabular}{@{}llccccr@{}}
\toprule
 & & \multicolumn{2}{c}{Coverage of the $95\%$ interval}
   & \multicolumn{2}{c}{KS $p$-value} & \\
\cmidrule(lr){3-4}\cmidrule(lr){5-6}
$h$ & $N$ & for $\theta$ & for $\alpha^{2}$ & $t_{N-1}$ & $\chi^{2}_{N-1}$
    & $\widehat\rho$\\
\midrule
$1/252$ & 30  & 0.9522 & 0.9513 & 0.719 & 0.024 & \phantom{$-$}0.0012\\
        & 100 & 0.9539 & 0.9488 & 0.807 & 0.379 & $-$0.0073\\
        & 200 & 0.9495 & 0.9508 & 0.212 & 0.181 & $-$0.0013\\
        & 300 & 0.9515 & 0.9479 & 0.807 & 0.544 & \phantom{$-$}0.0009\\
        & 500 & 0.9492 & 0.9486 & 0.991 & 0.777 & $-$0.0075\\
\addlinespace
$1/12$  & 30  & 0.9495 & 0.9489 & 0.764 & 0.558 & $-$0.0080\\
        & 100 & 0.9493 & 0.9496 & 0.087 & 0.548 & \phantom{$-$}0.0098\\
        & 200 & 0.9478 & 0.9523 & 0.136 & 0.757 & $-$0.0052\\
        & 300 & 0.9475 & 0.9497 & 0.146 & 0.633 & \phantom{$-$}0.0005\\
        & 500 & 0.9478 & 0.9498 & 0.394 & 0.097 & \phantom{$-$}0.0019\\
\bottomrule
\end{tabular}
\end{table}

\subsection{Graphical analysis}

Figure \ref{fi:paths} illustrates the model. Panel (a) shows six independent trajectories
of \eqref{eq:model} for Scenario 2 with $N=500$ and $h=1/252$, together with the trend
$\theta t$; the dispersion of the paths around the trend over a window of length
$t_N\approx1.98$ is a graphical expression of the variance bound \eqref{eq:varbound} and
explains why the drift is the hard parameter in this model. Panel (b) displays the three
fractional components of one such trajectory and shows how the roughness decreases with
$H_r$.

Figures \ref{fi:hist} and \ref{fi:qq} assess Theorem \ref{th:exactlaw} graphically. In
Figure \ref{fi:hist} the histograms of $\widehat\theta$ and $\widehat\alpha^{\,2}$ are
superimposed not on a fitted Gaussian density but on the \emph{exact} densities predicted
by the theorem, namely $\Norm(\theta,\alpha^{2}/d_N)$ and the density of
$(\alpha^{2}/N)\chi^{2}_{N-1}$. The agreement is complete at $N=500$ and equally complete
at $N=30$, where the pronounced right skewness of $\widehat\alpha^{\,2}$, which would
appear as a failure of normality, is precisely the skewness of the $\chi^{2}_{29}$
law. Figure \ref{fi:qq} makes the same point through quantile--quantile plots against the
exact reference laws: the four plots are straight lines, at $N=30$ as much as at $N=500$.

\subsection{Sensitivity to the Hurst vector}\label{sse:misspec}

Every exactness statement of Section \ref{se:mle} is conditional on $\GH$ being the true
covariance, that is, on the Hurst vector being known. In applications $H$ is postulated
or estimated, and it is legitimate to ask what survives when the matrix used for
inference is built from a wrong vector. We therefore generate data from
$H=(0.65,0.75,0.85)$ and compute the estimators and the intervals from
$H+\delta\bone$, with $\delta$ ranging over $\pm0.02$, $\pm0.05$ and $\pm0.10$;
Scenario 2, $h=1/252$, $N\in\{100,500\}$ and $20\,000$ replications throughout.
Table \ref{ta:misspec} reports the outcome.

The whole effect is carried by two scalars. Write $\Gamma_{\!a}$ for the assumed
covariance and $\GH$ for the true one, and set
\begin{equation}\label{eq:bq}
b_N(\delta)=\frac1N\tr\bigl(\Gamma_{\!a}^{-1}\GH\bigr),\qquad
q_N(\delta)=\frac{\bT^{\top}\Gamma_{\!a}^{-1}\GH\Gamma_{\!a}^{-1}\bT}
                 {\bT^{\top}\Gamma_{\!a}^{-1}\bT}.
\end{equation}
Both equal one when $\delta=0$. A direct computation gives
$\E[\widehat\alpha^{\,2}]=\alpha^{2}\{Nb_N(\delta)-q_N(\delta)\}/N$, and replacing the
quadratic form in the denominator of the studentised statistic by its mean yields the
approximation
\begin{equation}\label{eq:covapprox}
\text{coverage of \eqref{eq:ci-theta}}\;\approx\;
2F_{N-1}\Bigl(t_{N-1,\,1-\beta/2}\sqrt{b_N(\delta)/q_N(\delta)}\Bigr)-1,
\end{equation}
$F_{N-1}$ denoting the $t_{N-1}$ distribution function. Crude as it is,
\eqref{eq:covapprox} reproduces every entry of the simulated coverage column of
Table \ref{ta:misspec} to within $1.6$ Monte Carlo standard errors.

The two parameters again behave in opposite ways. For the scale, $b_N(\delta)$ is
essentially free of $N$: it equals $0.781$ at $N=100$ and $0.780$ at $N=500$ for
$\delta=-0.02$, and the simulated ratio $\E[\widehat\alpha^{\,2}]/\alpha^{2}$ follows it
exactly. The interval \eqref{eq:ci-alpha}, on the other hand, has width of order
$N^{-1/2}$, so a bias that does not vanish is eventually excluded from it and the
coverage collapses: already at $\delta=\pm0.02$ it has fallen to $0.62$ and $0.53$ at
$N=100$, and to about $0.02$ at $N=500$. The reason is identifiability rather than
numerical instability. Since $c(k)$ aggregates terms $h^{2H_r}$, shifting every $H_r$ by
$\delta$ multiplies the increment covariance by approximately $h^{2\delta}$, so that
$b_N(\delta)\approx h^{-2\delta}$, which gives $0.80$ and $1.25$ at $\delta=\mp0.02$
against the computed $0.78$ and $1.29$. A perturbation of the Hurst vector is in this
sense almost indistinguishable from a rescaling of $\alpha^{2}$: $\widehat\alpha^{\,2}$ should be read as an estimate of
$\alpha^{2}b_N(\delta)$, that is, of the scale relative to the postulated $H$.

For the drift the picture is far better. The degradation is gradual and, more usefully,
signed: the interval is too short when $b_N(\delta)<1$ and too long when
$b_N(\delta)>1$, so on a fine grid, where $h<1$ and $b_N(\delta)\approx h^{-2\delta}$,
understating the Hurst parameters shortens it and overstating them lengthens it. At
$\delta=-0.02$ the coverage of the nominal $95\%$ interval is $0.924$ at $N=100$ and
$0.913$ at $N=500$; at $\delta=+0.02$ it is $0.970$ and $0.975$. An error of five
hundredths, which is the order of the standard error of a Hurst estimator at these sample
sizes, still leaves $0.84$ on the pessimistic side and $0.99$ on the optimistic one. A
practical rule follows: when $H$ is uncertain, erring towards a larger $b_N(\delta)$,
which for $h<1$ means rounding the Hurst estimates upwards, yields a conservative
interval for $\theta$, at a cost in length that Table \ref{ta:misspec} quantifies.

The conclusion is that exactness is a statement made conditionally on $H$, and that the
two parameters inherit that conditioning very differently. Inference about the drift is
robust to small errors in the Hurst vector and errs in a controllable direction;
inference about the scale is not, because $\alpha^{2}$ and $H$ are nearly confounded. A
joint treatment of $(\theta,\alpha^{2},H)$, in which this confounding would be resolved
rather than assumed away, is the natural continuation of the present work.

\begin{table}[htbp]
\caption{Sensitivity to a misspecified Hurst vector. Data are generated with
$H=(0.65,0.75,0.85)$; the estimators and the nominal $95\%$ intervals are computed with
$H+\delta\bone$. Scenario 2 ($\theta=1.5$, $\alpha=1$), $h=1/252$, $20\,000$
replications; the Monte Carlo standard error of a coverage entry is at most $0.0033$.
The column ``approx.'' is the right-hand side of \eqref{eq:covapprox}.}
\label{ta:misspec}
\small
\begin{tabular}{@{}rrcccc@{}}
\toprule
 & & \multicolumn{2}{c}{Coverage for $\theta$}
   & Coverage & \\
\cmidrule(lr){3-4}
$N$ & $\delta$ & simulated & approx. & for $\alpha^{2}$
    & $\E[\widehat\alpha^{\,2}]/\alpha^{2}$\\
\midrule
100 & $-0.10$ & 0.7664 & 0.7710 & 0.0000 & 0.3004\\
    & $-0.05$ & 0.8703 & 0.8723 & 0.0095 & 0.5368\\
    & $-0.02$ & 0.9237 & 0.9230 & 0.6197 & 0.7724\\
    & $\phantom{-}0.00$ & 0.9499 & 0.9500 & 0.9487 & 0.9913\\
    & $+0.02$ & 0.9701 & 0.9706 & 0.5259 & 1.2800\\
    & $+0.05$ & 0.9898 & 0.9897 & 0.0057 & 1.9019\\
    & $+0.10$ & 0.9996 & 0.9994 & 0.0000 & 3.8306\\
\addlinespace
500 & $-0.10$ & 0.6908 & 0.6911 & 0.0000 & 0.3050\\
    & $-0.05$ & 0.8407 & 0.8385 & 0.0000 & 0.5422\\
    & $-0.02$ & 0.9131 & 0.9126 & 0.0241 & 0.7784\\
    & $\phantom{-}0.00$ & 0.9487 & 0.9500 & 0.9518 & 0.9980\\
    & $+0.02$ & 0.9750 & 0.9758 & 0.0200 & 1.2875\\
    & $+0.05$ & 0.9940 & 0.9948 & 0.0000 & 1.9113\\
    & $+0.10$ & 1.0000 & 1.0000 & 0.0000 & 3.8458\\
\bottomrule
\end{tabular}
\end{table}

\section{Discussion and concluding remarks}\label{se:disc}

For a linear drift observed together with a superposition of independent fractional
Brownian motions sharing a common scale, likelihood inference is exact. The maximum
likelihood estimators of the drift and of the scale are available in closed form; their
joint law is known at every sample size, with $\widehat\theta$ Gaussian,
$N\widehat\alpha^{\,2}/\alpha^{2}$ chi-square with $N-1$ degrees of freedom and the two
independent; and confidence intervals and tests of exact level follow, together with
complete sufficiency, minimum variance unbiasedness and attainment of the Cram\'er--Rao
bound. The asymptotic theory, mean square and strong consistency and asymptotic
normality, is then a consequence rather than a substitute.

Several features of the analysis deserve emphasis. First, the model extends the classical
drifted fractional Brownian motion framework to an arbitrary finite number of independent
components, which enlarges the range of accessible covariance structures while preserving
Gaussianity and hence tractability of the likelihood; the single-component, mixed and
double-fractional models arise as special cases (Remark \ref{re:special}).

Second, the entire multi-mixed structure is compressed into the single scalar
$d_N=\bT^{\top}\GH^{-1}\bT$, which by Proposition \ref{pr:increment} is
$h^{2}\bone^{\top}C_N^{-1}\bone$ for the Toeplitz increment covariance $C_N$. Once the
noise covariance is known up to a positive factor, whitening reduces the model to a
Gaussian linear model with one regressor (Remark \ref{re:glm}), and the number of
components, the Hurst vector and the sampling step are all absorbed into $d_N$; this is
what delivers exact distributions and exact optimality for arbitrary $m$ and arbitrary
Hurst vectors. The same observation delimits the theory: with component-specific scales
the covariance is no longer a known matrix times an unknown constant, and exactness is
lost.

Third, the two parameters behave very differently. Inference about the scale is free of
the model: after whitening by $\GH^{-1/2}$ the entire multi-scale dependence structure
disappears from the sampling distribution of $\widehat\alpha^{\,2}$, which depends on $N$
alone, so the chi-square interval has exact level whatever the dependence structure.
Inference about the drift is not free. The explicit bound \eqref{eq:varbound} shows that
$\Var(\widehat\theta)$ is controlled by $\sum_rt_N^{2H_r-2}$, so that the component with
the largest Hurst parameter dominates, and by Remark \ref{re:sharp} the resulting order
$(Nh)^{2\Hmax-2}$ is exact and not merely an upper bound. What can be learned about
$\theta$ is therefore decided by the length of the observation window and by the
strongest long-range dependence present in the superposition; refining the mesh on a
fixed window does not help. This is the precise sense in which stronger long-range
dependence reduces the information available about the drift, and it is clearly visible
in the simulations.

Fourth, the sequence of drift estimators indexed by the sample size turns out to be, in
law, a Brownian motion run along its own decreasing variance scale
(Corollary \ref{co:tcbm}). This is a somewhat unexpected structure in a strongly
dependent data setting and it gives a transparent, moment-free proof of strong
consistency.

Fifth, all of this holds given the Hurst vector, and Section \ref{sse:misspec} measures
what a wrong one costs. The asymmetry between the two parameters appears there from the
other side: the Student interval for $\theta$ tolerates a small error in $H$ and loses
level in a direction one can predict, whereas the chi-square interval for $\alpha^{2}$
does not, since shifting every $H_r$ by $\delta$ multiplies the increment covariance by
approximately $h^{2\delta}$ and is therefore almost indistinguishable from a rescaling of
$\alpha^{2}$.

Several directions remain open. The most natural is to treat $H_1,\dots,H_m$ as unknown
and estimate them jointly with $\theta$ and $\alpha^{2}$, which would also resolve the
confounding just described; since $\GH$ depends nonlinearly on $H$, both the numerical
optimisation and the asymptotic analysis become substantially harder, and even
identifiability of the individual $H_r$ from a single trajectory requires care. A second
direction is to relax the common-scale restriction and allow component-specific scales
$\alpha_1,\dots,\alpha_m$. Here Remark \ref{re:glm} indicates what is at stake: the
covariance matrix $\sum_r\alpha_r^{2}\Gamma_{H_r}$ is then no longer a scalar multiple of
a known matrix, the profile likelihood in the scale parameters no longer admits a closed
form, and one should expect the individual $\alpha_r$ to be weakly identified when the
corresponding Hurst parameters are close, since $\Gamma_{H}$ depends continuously on $H$.
Quantifying that loss of information, for instance through the Fisher information matrix
of $(\alpha_1^{2},\dots,\alpha_m^{2})$, is a natural next step. A third direction is to
relax the equidistant sampling assumption: the likelihood framework of
Theorem \ref{th:loglik} extends verbatim by modifying $\GH$, but
Proposition \ref{pr:increment} and the explicit bounds of Proposition \ref{pr:varbound}
rely on the regular grid.

\begin{thebibliography}{99}

\bibitem{Ade74}
R.~K. Adenstedt,
\emph{On large-sample estimation for the mean of a stationary random sequence},
Ann.\ Statist.\ \textbf{2} (1974), no.~6, 1095--1107.

\bibitem{AM21}
S.~Alajmi and E.~Mliki,
\emph{Mixed generalized fractional Brownian motion},
J.\ Stoch.\ Anal.\ \textbf{2} (2021), no.~2, Article~2.

\bibitem{AG07}
S.~Asmussen and P.~W. Glynn,
\emph{Stochastic Simulation: Algorithms and Analysis},
Stochastic Modelling and Applied Probability, vol.~57, Springer, New York, 2007.

\bibitem{HN25}
M.~R. Haddadi and H.~Nasrollahi,
\emph{Option pricing under non-normal distribution in mixed of Gram--Charlier model and
fractional models},
J.\ Math.\ Model.\ Finance \textbf{5} (2025), no.~1, 47--62.

\bibitem{Hig01}
D.~J. Higham,
\emph{An algorithmic introduction to numerical simulation of stochastic differential
equations},
SIAM Rev.\ \textbf{43} (2001), no.~3, 525--546.

\bibitem{HNXZ11}
Y.~Hu, D.~Nualart, W.~Xiao, and W.~Zhang,
\emph{Exact maximum likelihood estimator for drift fractional Brownian motion at discrete
observation},
Acta Math.\ Sci.\ Ser.\ B \textbf{31} (2011), no.~5, 1851--1859.

\bibitem{LC98}
E.~L. Lehmann and G.~Casella,
\emph{Theory of Point Estimation}, 2nd ed.,
Springer Texts in Statistics, Springer, New York, 1998.

\bibitem{LR05}
E.~L. Lehmann and J.~P. Romano,
\emph{Testing Statistical Hypotheses}, 3rd ed.,
Springer Texts in Statistics, Springer, New York, 2005.

\bibitem{KL15}
N.~Kuang and B.~Liu,
\emph{Parameter estimations for the sub-fractional Brownian motion with drift at discrete
observation},
Braz.\ J.\ Probab.\ Stat.\ \textbf{29} (2015), no.~4, 778--789.

\bibitem{MS23}
H.~Maleki Almani and T.~Sottinen,
\emph{Multi-mixed fractional Brownian motions and Ornstein--Uhlenbeck processes},
Mod.\ Stoch.\ Theory Appl.\ \textbf{10} (2023), no.~4, 343--366.

\bibitem{MVN68}
B.~B. Mandelbrot and J.~W. Van Ness,
\emph{Fractional Brownian motions, fractional noises and applications},
SIAM Rev.\ \textbf{10} (1968), no.~4, 422--437.

\bibitem{Mis08}
Y.~Mishura,
\emph{Stochastic Calculus for Fractional Brownian Motion and Related Processes},
Lecture Notes in Mathematics, vol.~1929, Springer, Berlin, 2008.

\bibitem{MRS17}
Y.~Mishura, K.~Ralchenko, and S.~Shklyar,
\emph{Maximum likelihood drift estimation for Gaussian process with stationary
increments},
Austrian J.\ Stat.\ \textbf{46} (2017), no.~3--4, 67--78.

\bibitem{MV15}
Y.~Mishura and I.~Voronov,
\emph{Construction of maximum likelihood estimator in the mixed fractional--fractional
Brownian motion model with double long-range dependence},
Mod.\ Stoch.\ Theory Appl.\ \textbf{2} (2015), no.~2, 147--164.

\bibitem{NP12}
I.~Nourdin and G.~Peccati,
\emph{Normal Approximations with Malliavin Calculus: From Stein's Method to Universality},
Cambridge Tracts in Mathematics, vol.~192, Cambridge University Press, Cambridge, 2012.

\bibitem{NOL08}
D.~Nualart and S.~Ortiz-Latorre,
\emph{Central limit theorems for multiple stochastic integrals and Malliavin calculus},
Stochastic Process.\ Appl.\ \textbf{118} (2008), no.~4, 614--628.

\bibitem{RY23}
K.~Ralchenko and M.~Yakovliev,
\emph{Asymptotic normality of parameter estimators for mixed fractional Brownian motion
with trend},
Austrian J.\ Stat.\ \textbf{52} (2023), 127--148.

\bibitem{RY24}
K.~Ralchenko and M.~Yakovliev,
\emph{Parameter estimation for fractional mixed fractional Brownian motion based on
discrete observations},
Mod.\ Stoch.\ Theory Appl.\ \textbf{11} (2024), no.~1, 1--29.

\bibitem{XZZ11}
W.~Xiao, W.~Zhang, and X.~Zhang,
\emph{Maximum-likelihood estimators in the mixed fractional Brownian motion},
Statistics \textbf{45} (2011), no.~1, 73--85.

\end{thebibliography}

\end{document}
